% Luc Destombe --- licence CC BY-NC-SA 4.0
% https://creativecommons.org/licenses/by-nc-sa/4.0/deed.fr
% Fichier autonome : compiler avec pdflatex, deux fois (signets, nombre de pages).
\documentclass[11pt,a4paper]{article}
\makeatletter
\newif\iflycee@unecolonne
\lycee@unecolonnetrue
\newif\iflycee@palatino
\lycee@palatinotrue

\RequirePackage[utf8]{inputenc}
\RequirePackage[T1]{fontenc}
\RequirePackage[french]{babel}
\RequirePackage{etoolbox}

\newcommand{\ShorthandsOff}{\@ifundefined{shorthandoff}{}{\shorthandoff{;:!?}}}
\newcommand{\ShorthandsOn}{\@ifundefined{shorthandon}{}{\shorthandon{;:!?}}}
\ShorthandsOff
\AtBeginDocument{\ShorthandsOn}
\AtBeginEnvironment{tikzpicture}{\ShorthandsOff}

\RequirePackage{geometry}
\RequirePackage{fancyhdr}
\RequirePackage{lastpage}
\RequirePackage{multicol}
\RequirePackage{array}
\RequirePackage{enumitem}
\RequirePackage{xcolor}
\RequirePackage{needspace}

\RequirePackage{amsmath,amssymb}
\RequirePackage{mathpazo}
\DeclareMathSizes{10.95}{10.95}{8}{5.5}
\begingroup\catcode`\_=\active \catcode`\^=\active
\gdef_#1{\sb{\scriptscriptstyle #1}}
\gdef^#1{\sp{\scriptscriptstyle\lycee@moinsepais #1}}
\endgroup
\newcommand\lycee@moins{\mathbin{%
  \setbox\z@\hbox{$\m@th\scriptscriptstyle\lycee@moinsorig$}%
  \dimen@=.55\wd\z@ \advance\dimen@-.5pt
  \hbox to.75\wd\z@{\hss$\m@th\scriptscriptstyle\vcenter{\hbox{%
    \tikz\draw[line width=.5pt,line cap=round](0,0)--(\the\dimen@,0);}}$\hss}}}
\begingroup\lccode`\~=`\- \lowercase{\endgroup\let~}\lycee@moins
\newcommand\lycee@moinsepais{\mathcode`\-="8000 }
\AtBeginDocument{\mathchardef\lycee@moinsorig=\mathcode`\-\relax}
\AtBeginDocument{\catcode`\_=12 \mathcode`\_="8000
  \catcode`\^=12 \mathcode`\^="8000 }
\newcommand\lycee@exposants{%
  \fontdimen13\textfont2=.48\fontdimen6\textfont2
  \fontdimen14\textfont2=.48\fontdimen6\textfont2
  \fontdimen15\textfont2=.48\fontdimen6\textfont2
  \fontdimen13\scriptfont2=.48\fontdimen6\scriptfont2
  \fontdimen14\scriptfont2=.48\fontdimen6\scriptfont2
  \fontdimen15\scriptfont2=.48\fontdimen6\scriptfont2}
\every@math@size\expandafter{\the\every@math@size\lycee@exposants}
\newlength{\retraitcalculs}
\setlength{\retraitcalculs}{2em}

\RequirePackage{tikz}
\usetikzlibrary{arrows.meta,calc,babel,shapes.misc}
\RequirePackage[most]{tcolorbox}
\tcbuselibrary{skins,breakable}

\definecolor{coulDef}{HTML}{1F4E79}
\definecolor{coulProp}{HTML}{7030A0}
\definecolor{coulRem}{HTML}{C55A11}
\definecolor{coulEx}{HTML}{2E7D32}
\definecolor{coulExo}{HTML}{B03A2E}
\definecolor{coulPart}{HTML}{1F4E79}
\definecolor{coulMeth}{HTML}{1B6E4F}
\definecolor{coulCourbe}{HTML}{0B5ED7}
\definecolor{coulCourbeB}{HTML}{C00000}
\definecolor{coulCourbeC}{HTML}{2E7D32}
\definecolor{coulGrille}{HTML}{8C8C8C}
\definecolor{coulRep}{HTML}{1B6E4F}
\definecolor{coulBar}{HTML}{2E7D32}

\newcommand{\R}{\mathbb{R}}
\newcommand{\Rs}{\mathbb{R}^{*}}
\renewcommand{\le}{\leqslant}
\renewcommand{\ge}{\geqslant}
\renewcommand{\approx}{\simeq}
\newcommand{\vect}[1]{\overrightarrow{#1}}
\newcommand{\norme}[1]{\left\lVert #1 \right\rVert}
\newcommand{\intervalle}[2]{\left[#1\,;#2\right]}
\RequirePackage{cancel}
\renewcommand{\CancelColor}{\color{coulExo}}
\newcommand{\fois}[1]{\times{\color{coulExo}#1}}
\newcommand{\barre}[1]{\cancel{#1}}

\tikzset{grille/.style={coulGrille,line width=0.4pt}}
\newcommand{\quadrillage}[4]{\draw[grille,step=1] (#1,#2) grid (#3,#4);}
\tikzset{
  croix/.style={cross out,draw,line width=0.6pt,inner sep=0pt,minimum size=4pt},
  croixdroite/.style={croix,rotate=45,line width=0.8pt},
}
\newcommand{\pt}[3]{\path (#1) node[croix]{} node[#2,font=\small,inner sep=2.5pt]{$#3$};}

\newcommand{\lycee@repere}[4]{%
  \draw[grille,step=1] (#1,#2) grid (#3,#4);
  \draw[-{Stealth[length=2mm]},thick] (#1,0)--({#3+0.4},0) node[below left=-1pt]{$x$};
  \draw[-{Stealth[length=2mm]},thick] (0,#2)--(0,{#4+0.4}) node[below left=-1pt]{$y$};
  \node[below left,font=\scriptsize,inner sep=1.5pt] at (0,0) {$0$};}
\newcommand{\lycee@gradx}[1]{\foreach \k/\l in {#1}{%
  \draw (\k,0.07)--(\k,-0.07) node[below,font=\scriptsize,inner sep=1.5pt]{$\l$};}}
\newcommand{\lycee@grady}[1]{\foreach \k/\l in {#1}{%
  \draw (0.07,\k)--(-0.07,\k) node[left,font=\scriptsize,inner sep=1.5pt]{$\l$};}}
\newcommand{\lycee@intff}[2]{\left[#1\,;#2\right]}
\newcommand{\lycee@intfo}[2]{\left[#1\,;#2\right[}
\newcommand{\lycee@intof}[2]{\left]#1\,;#2\right]}
\newcommand{\lycee@intoo}[2]{\left]#1\,;#2\right[}
\newcommand{\lycee@ens}[1]{\left\{#1\right\}}
\AtBeginDocument{%
  \providecommand{\repere}{\lycee@repere}%
  \providecommand{\gradx}{\lycee@gradx}%
  \providecommand{\grady}{\lycee@grady}%
  \providecommand{\Cf}{\mathcal{C}\sb{\scriptscriptstyle f}}%
  \providecommand{\Cg}{\mathcal{C}\sb{\scriptscriptstyle g}}%
  \providecommand{\intff}{\lycee@intff}%
  \providecommand{\intfo}{\lycee@intfo}%
  \providecommand{\intof}{\lycee@intof}%
  \providecommand{\intoo}{\lycee@intoo}%
  \providecommand{\ens}{\lycee@ens}}

\newlist{questions}{enumerate}{3}
\setlist[questions,1]{label=\arabic*\textdegree),leftmargin=*,itemsep=2pt,topsep=3pt}
\setlist[questions,2]{label=\alph*),leftmargin=*,itemsep=1pt,topsep=2pt}
\setlist[questions,3]{label=\roman*),leftmargin=*,itemsep=1pt,topsep=1pt}
\setlist[itemize]{label=\textbullet,leftmargin=*,itemsep=2pt,topsep=3pt}

\newcommand{\besoinplace}[1]{\par\penalty\@M\begingroup
  \dimen@=#1\relax\dimen@ii=\pagegoal\advance\dimen@ii-\pagetotal
  \ifdim\dimen@>\dimen@ii\ifdim\dimen@ii>\z@\vfil\fi\break\fi\endgroup}

\newcommand{\lycee@pied}{\thepage/\pageref{LastPage}}
\newcommand{\entete}[2]{%
  \pagestyle{fancy}\fancyhf{}%
  \fancyhead[L]{\footnotesize\color{coulPart}\textbf{#1}}%
  \fancyhead[R]{\footnotesize\color{coulPart}\textbf{#2}}%
  \fancyfoot[C]{\footnotesize\lycee@pied}%
  \renewcommand{\headrulewidth}{0.4pt}%
  \renewcommand{\headrule}{\hbox to\headwidth{%
    \color{coulPart!40}\leaders\hrule height \headrulewidth\hfill}}}

\RequirePackage{ccicons}
\newcommand{\lycee@licence}{{\scriptsize\color{gray}%
  \copyright\ Luc Destombe --- \ccbyncsa\ CC BY-NC-SA 4.0}}
\newif\iflycee@licence \lycee@licencetrue
\newcommand{\sanslicence}{\lycee@licencefalse}
\AtBeginDocument{\iflycee@licence\AddToHookNext{shipout/foreground}{%
  \put(0,0){\rlap{\kern\dimexpr1in+\hoffset+\oddsidemargin+\textwidth\relax
    \llap{\raisebox{-\dimexpr1in+\voffset+\topmargin+\headheight+\headsep
      +\textheight+\footskip\relax}{\lycee@licence}}}}}\fi}

\newcommand{\formulesserrees}{%
  \setlength{\abovedisplayskip}{3pt}\setlength{\belowdisplayskip}{3pt}%
  \setlength{\abovedisplayshortskip}{2pt}\setlength{\belowdisplayshortskip}{3pt}}

\setlength{\parindent}{0pt}

\RequirePackage[implicit=false,hidelinks,pdfencoding=auto,psdextra,bookmarksopen,
  bookmarksopenlevel=1,pdfcreator={},pdfproducer={}]{hyperref}
\RequirePackage{bookmark}
\hypersetup{pdfauthor={Luc Destombe},pdfkeywords={CC BY-NC-SA 4.0}}
\newcounter{lycee@signet}
\newcommand{\signet}[3][]{\stepcounter{lycee@signet}%
  \edef\lycee@dest{\if\relax\detokenize{#1}\relax
    signet.\arabic{lycee@signet}\else#1\fi}%
  \hypertarget{\lycee@dest}{}\bookmark[level=#2,dest=\lycee@dest]{#3}}
\pdfstringdefDisableCommands{%
  \def\R{R}\def\vect#1{#1}\def\Cf{Cf}\def\Cg{Cg}\def\fois#1{×#1}%
  \def\barre#1{#1}\def\intervalle#1#2{[#1;#2]}\def\intff#1#2{[#1;#2]}%
  \def\textsuperscript#1{#1}\def\dfrac#1#2{#1/#2}\def\frac#1#2{#1/#2}%
  \def\vec#1{#1⃗}}%
\RequirePackage[official]{eurosym}

\tcbset{exercice corrige/.style={
  enhanced,breakable,before upper=\raggedright,
  colback=white,colframe=coulExo,
  boxrule=0.8pt,arc=2pt,
  left=6pt,right=6pt,top=8pt,bottom=5pt,
  before skip=9pt,after skip=4pt,
  coltitle=white,fonttitle=\bfseries\small,
  attach boxed title to top left={xshift=8pt,yshift=-\tcboxedtitleheight/2},
  boxed title style={colback=coulExo,arc=2pt,boxrule=0pt}}}
\newcounter{exo}
\newtcolorbox{exercice}[1][]{exercice corrige,
  phantom={\signet[exercice-\arabic{exo}]{2}{Exercice \arabic{exo}}},
  title={Exercice \theexo\ifx\relax#1\relax\else{} --- #1\fi}}
\newenvironment{exo}[1][\relax]{\refstepcounter{exo}\begin{exercice}[#1]}{\end{exercice}}

\geometry{top=1.8cm,bottom=1cm,left=1cm,right=1cm,headheight=14pt,footskip=0.6cm}
\renewcommand{\lycee@pied}{\thepage}

\renewcommand{\theexo}{\ifnum\value{exo}<10 0\fi\arabic{exo}}
\tcbset{exercice corrige/.style={
  enhanced,breakable,before upper=\raggedright,
  colback=white,colframe=coulExo,
  boxrule=0.9pt,arc=2pt,
  left=8pt,right=8pt,top=8pt,bottom=6pt,
  before skip=8pt,after skip=6pt,
  coltitle=white,fonttitle=\bfseries,
  attach boxed title to top left={xshift=8pt,yshift=-\tcboxedtitleheight/2},
  boxed title style={colback=coulExo,arc=2pt,boxrule=0pt}}}

\newcommand{\entetefiche}[2]{%
  \begin{center}
    \begin{tcolorbox}[enhanced,width=0.92\linewidth,colback=white,colframe=coulPart,
        boxrule=1.4pt,arc=3pt,halign=center,top=6pt,bottom=6pt]
      {\large\bfseries\color{coulPart} CORRIGÉ --- FICHE D'EXERCICES}\\[3pt]
      {\normalsize\bfseries #1}\\[2pt]
      {\normalsize #2}
    \end{tcolorbox}
  \end{center}}

\newcounter{partie}
\newcommand{\partie}[1]{%
  \stepcounter{partie}%
  \besoinplace{8\baselineskip}\vspace{4pt}%
  \begin{tcolorbox}[enhanced,colback=coulPart!8,colframe=coulPart,boxrule=0pt,
      leftrule=3pt,arc=0pt,left=6pt,right=4pt,top=3pt,bottom=3pt,
      before skip=6pt,after skip=2pt]
    \bfseries\color{coulPart}\Roman{partie} --- #1\signet{1}{\Roman{partie} --- #1}
  \end{tcolorbox}}
\newcommand{\partiesansnum}[1]{%
  \besoinplace{8\baselineskip}\vspace{4pt}%
  \begin{tcolorbox}[enhanced,colback=coulPart!8,colframe=coulPart,boxrule=0pt,
      leftrule=3pt,arc=0pt,left=6pt,right=4pt,top=3pt,bottom=3pt,
      before skip=6pt,after skip=2pt]
    \bfseries\color{coulPart}#1\signet{1}{#1}
  \end{tcolorbox}}

\setlist[questions,1]{label=\arabic*\textdegree),leftmargin=*,itemsep=3pt,topsep=3pt}
\setlist[questions,2]{label=\alph*),leftmargin=*,itemsep=2pt,topsep=2pt}
\setlist[questions,3]{label=\roman*),leftmargin=*,itemsep=1pt,topsep=1pt}
\newlist{expr}{enumerate}{1}
\setlist[expr]{label={\Alph*\ $=$},leftmargin=*,itemsep=2pt,topsep=3pt}

\newcommand{\res}[1]{{\color{coulRep}#1}}
\newcommand{\rem}[1]{\par\smallskip{\small\itshape\color{black!60}$\rightarrow$ #1}}
\newenvironment{calculs}{\par\smallskip\noindent$\displaystyle\begin{aligned}[t]}%
  {\end{aligned}$\par\smallskip}
\newcommand{\justif}[1]{\quad\text{\small\itshape\color{black!60}#1}}
\newcommand{\cas}[1]{\par\smallskip\noindent\textbf{\color{coulExo}#1}\nobreak}
\newcommand{\explique}[2]{\par\smallskip\noindent
  \begin{minipage}[t]{0.6\linewidth}\raggedright #1\end{minipage}\hfill
  \begin{minipage}[t]{0.37\linewidth}\raggedright\leavevmode
    {\small\itshape\color{black!60}#2}\end{minipage}\par}

\setlength{\parskip}{2pt}

\makeatother
\entete{Chapitre 3 --- Vecteurs --- corrigé}{Seconde}
\setlength{\parindent}{0pt}
\setlength{\jot}{5pt}

\let\casligne\cas
\renewcommand{\cas}[1]{\casligne{#1}\par\nopagebreak}

\tikzset{
  vecteur/.style={-{Stealth[length=2.2mm]},line width=0.9pt,coulCourbe},
  vecteurB/.style={-{Stealth[length=2.2mm]},line width=0.9pt,coulCourbeB},
  vecteurC/.style={-{Stealth[length=2.2mm]},line width=0.9pt,coulCourbeC},
  cherche/.style={-{Stealth[length=2.4mm]},line width=1.1pt,black},
  aide/.style={dashed,black!50,line width=0.5pt},
  fig/.style={line width=0.7pt,black!75},
  figC/.style={line width=0.9pt,coulCourbeC},
  nom/.style={midway,font=\footnotesize,inner sep=1.5pt},
}

\newcommand{\cotefig}[3][0.32]{\par\vspace{-0.55\baselineskip}\noindent
  \begin{minipage}[t]{\dimexpr0.97\linewidth-#1\linewidth\relax}\raggedright #2\end{minipage}\hfill
  \begin{minipage}[t]{#1\linewidth}\vspace{0pt}\centering #3\end{minipage}\par}

\newcommand{\figKLM}{%
  \quadrillage{-0.4}{-0.4}{8.4}{2.4}
  \draw[fig] (0,0)--(6,0) (2,2)--(8,2) (0,0)--(2,2) (3,0)--(5,2) (6,0)--(8,2);
  \pt{0,0}{below}{K} \pt{3,0}{below}{L} \pt{6,0}{below}{M}
  \pt{2,2}{above}{N} \pt{5,2}{above}{P} \pt{8,2}{above}{Q}}
\newcommand{\triangleABC}{%
  \draw[fig] (0,0)--(4,0)--(1,3)--cycle;
  \pt{0,0}{below left}{A} \pt{4,0}{below right}{B} \pt{1,3}{above right}{C}}
\newcommand{\triangleDEF}{%
  \draw[fig] (0,0)--(4,0)--(1,3)--cycle;
  \pt{0,0}{below left}{D} \pt{4,0}{below right}{E} \pt{1,3}{above}{F}}
\newcommand{\figABCD}{%
  \quadrillage{-0.4}{-0.4}{8.4}{5.4}
  \pt{0,1}{below left}{A} \pt{4,1}{below}{B} \pt{6,3}{right}{C} \pt{1,4}{above left}{D}}
\newcommand{\figABCDbis}{%
  \quadrillage{-0.4}{-2.4}{8.4}{4.4}
  \pt{0,1}{left}{A} \pt{4,0}{below}{B} \pt{5,3}{above right}{C} \pt{1,3}{above}{D}}

\begin{document}

\entetefiche{Chapitre 3 --- Vecteurs}{Seconde}

{\small Sur un quadrillage, on lit chaque vecteur comme un \textbf{déplacement} : tant
de carreaux à droite ou à gauche, tant vers le haut ou vers le bas. Sur les figures, le
vecteur donné est en couleur ; un représentant reporté garde sa couleur et son nom ; le
vecteur cherché est en noir.\par}

\partie{Caractéristiques d'un vecteur}

\begin{exo}[Égaux ou non]
\explique{Pour chaque égalité, on trace les deux flèches, de la première lettre vers la
seconde, puis on compare leurs déplacements.}{méthode 1 : même déplacement, vecteurs
égaux}
\cas{a)}
\cotefig{%
\explique{$\vect{KL}$ : $3$ carreaux à droite.}{de $K$ vers $L$}
\explique{$\vect{NP}$ : $3$ carreaux à droite.}{de $N$ vers $P$}
\explique{Même déplacement, donc $\vect{KL}\mathrel{\res{=}}\vect{NP}$.}{même direction, même sens,
même norme}}
{\begin{tikzpicture}[scale=0.42]\figKLM
  \draw[vecteur] (0,0)--(3,0) node[nom,above]{$\vect{KL}$};
  \draw[vecteurB] (2,2)--(5,2) node[nom,below]{$\vect{NP}$};
\end{tikzpicture}}
\cas{b)}
\cotefig{%
\explique{$\vect{KN}$ : $2$ carreaux à droite et $2$ vers le haut.}{de $K$ vers $N$}
\explique{$\vect{LP}$ : $2$ carreaux à droite et $2$ vers le haut.}{de $L$ vers $P$}
\explique{Donc $\vect{KN}\mathrel{\res{=}}\vect{LP}$.}{même déplacement}}
{\begin{tikzpicture}[scale=0.42]\figKLM
  \draw[vecteur] (0,0)--(2,2) node[nom,above left]{$\vect{KN}$};
  \draw[vecteurB] (3,0)--(5,2) node[nom,right]{$\vect{LP}$};
\end{tikzpicture}}
\cas{c)}
\cotefig{%
\explique{$\vect{KP}$ : $5$ carreaux à droite et $2$ vers le haut.}{de $K$ vers $P$}
\explique{$\vect{LQ}$ : $5$ carreaux à droite et $2$ vers le haut.}{de $L$ vers $Q$}
\explique{Donc $\vect{KP}\mathrel{\res{=}}\vect{LQ}$.}{même déplacement}}
{\begin{tikzpicture}[scale=0.42]\figKLM
  \draw[vecteur] (0,0)--(5,2) node[nom,above left]{$\vect{KP}$};
  \draw[vecteurB] (3,0)--(8,2) node[nom,below right]{$\vect{LQ}$};
\end{tikzpicture}}
\cas{d)}
\cotefig{%
\explique{$\vect{NL}$ : $1$ carreau à droite et $2$ vers le bas.}{de $N$ vers $L$}
\explique{$\vect{KN}$ : $2$ carreaux à droite et $2$ vers le haut.}{de $K$ vers $N$}
\explique{Les déplacements sont différents, donc $\vect{NL}\mathrel{\res{\neq}}\vect{KN}$.}{les
directions ne sont même pas les mêmes}}
{\begin{tikzpicture}[scale=0.42]\figKLM
  \draw[vecteur] (2,2)--(3,0) node[nom,right]{$\vect{NL}$};
  \draw[vecteurB] (0,0)--(2,2) node[nom,above left]{$\vect{KN}$};
\end{tikzpicture}}
\cas{e)}
\cotefig{%
\explique{$\vect{KM}$ : $6$ carreaux à droite.}{de $K$ vers $M$}
\explique{$\vect{NQ}$ : $6$ carreaux à droite.}{de $N$ vers $Q$}
\explique{Donc $\vect{KM}\mathrel{\res{=}}\vect{NQ}$.}{même déplacement}}
{\begin{tikzpicture}[scale=0.42]\figKLM
  \draw[vecteur] (0,0)--(6,0) node[nom,above]{$\vect{KM}$};
  \draw[vecteurB] (2,2)--(8,2) node[nom,below]{$\vect{NQ}$};
\end{tikzpicture}}
\cas{f)}
\cotefig{%
\explique{$\vect{PQ}$ : $3$ carreaux à droite.}{de $P$ vers $Q$}
\explique{$\vect{PN}$ : $3$ carreaux à gauche.}{de $P$ vers $N$}
\explique{Donc $\vect{PQ}\mathrel{\res{\neq}}\vect{PN}$.}{les sens sont contraires}}
{\begin{tikzpicture}[scale=0.42]\figKLM
  \draw[vecteur] (5,2)--(8,2) node[nom,below]{$\vect{PQ}$};
  \draw[vecteurB] (5,2)--(2,2) node[nom,below]{$\vect{PN}$};
\end{tikzpicture}}
\rem{f) Même direction, même norme, mais sens contraires : $\vect{PN}$ et $\vect{PQ}$ sont
opposés, $\vect{PN}=-\vect{PQ}$. Deux flèches de même longueur sur une même droite ne
suffisent pas : il faut regarder où est la pointe.}
\end{exo}

\begin{exo}[Lire sur la figure]
\explique{On garde la figure de l'exercice 1 et on lit chaque déplacement.}{$K$, $L$, $M$
en bas ; $N$, $P$, $Q$ en haut, $3$ carreaux entre deux points voisins}
\cas{1)}
\explique{$\vect{KN}$ fait $2$ carreaux à droite et $2$ vers le haut ; $\vect{LP}$ et
$\vect{MQ}$ aussi. Donc $\vect{KN}=\res{\vect{LP}}$ (ou $\vect{MQ}$).}{les côtés obliques
des parallélogrammes}
\explique{$\vect{LM}$ fait $3$ carreaux à droite ; $\vect{KL}$, $\vect{NP}$ et $\vect{PQ}$
aussi. Donc $\vect{LM}=\res{\vect{KL}}$ (ou $\vect{NP}$, $\vect{PQ}$).}{un seul exemple
suffit}
\cas{2)}
\explique{Un vecteur opposé à $\vect{KN}$ fait $2$ carreaux à gauche et $2$ vers le bas :
$\res{\vect{NK}}$, ou $\vect{PL}$, ou $\vect{QM}$.}{on échange les deux lettres : le sens
change}
\explique{Un vecteur opposé à $\vect{KL}$ fait $3$ carreaux à gauche : $\res{\vect{LK}}$,
ou $\vect{ML}$, $\vect{PN}$, $\vect{QP}$.}{même direction, même norme, sens contraire}
\cas{3)}
\explique{Même direction et même sens, mais pas égaux : $\res{\vect{KL}}$ et
$\res{\vect{KM}}$.}{tous deux vers la droite, mais $3$ et $6$ carreaux : normes
différentes}
\explique{Même norme, mais pas égaux : $\res{\vect{KL}}$ et $\res{\vect{PN}}$.}{tous deux
de $3$ carreaux, mais de sens contraires}
\explique{Même direction et sens contraires, mais pas opposés : $\res{\vect{KM}}$ et
$\res{\vect{LK}}$.}{$6$ carreaux à droite, $3$ à gauche : normes différentes}
\rem{Deux vecteurs égaux ont les trois caractéristiques en commun ; il suffit qu'une seule
diffère pour qu'ils ne le soient plus. Ici, d'autres réponses conviennent.}
\end{exo}

\begin{exo}[Construire un représentant]
\explique{On trace d'abord $\vec{u}$ : une flèche de $3$~cm. Puis on construit chaque
$\vec{v}$ en partant de ses trois caractéristiques.}{figures réduites de moitié}
\cas{1)}
\cotefig{%
\explique{On trace une flèche parallèle à $\vec{u}$, pointée du même côté, et deux fois
plus longue : $6$~cm.}{même direction, même sens, norme $6$~cm}}
{\begin{tikzpicture}[scale=0.5]
  \draw[vecteur] (0,1)--(3,1) node[nom,above]{$\vec{u}$};
  \draw[cherche] (0,0)--(6,0) node[nom,below]{$\vec{v}$};
\end{tikzpicture}}
\cas{2)}
\cotefig{%
\explique{On trace une flèche parallèle à $\vec{u}$, de $3$~cm, pointée de l'autre
côté.}{c'est un représentant de $-\vec{u}$}}
{\begin{tikzpicture}[scale=0.5]
  \draw[vecteur] (0,1)--(3,1) node[nom,above]{$\vec{u}$};
  \draw[cherche] (3,0)--(0,0) node[nom,below]{$\vec{v}$};
\end{tikzpicture}}
\cas{3)}
\cotefig{%
\explique{On trace une flèche parallèle à $\vec{u}$, pointée de l'autre côté, de
$1{,}5$~cm.}{sens contraire, deux fois plus courte}}
{\begin{tikzpicture}[scale=0.5]
  \draw[vecteur] (0,1)--(3,1) node[nom,above]{$\vec{u}$};
  \draw[cherche] (1.5,0)--(0,0) node[nom,below]{$\vec{v}$};
\end{tikzpicture}}
\cas{4)}
\cotefig{%
\explique{On trace une flèche de $3$~cm qui n'est pas parallèle à $\vec{u}$.}{toute autre
direction convient}}
{\begin{tikzpicture}[scale=0.5]
  \draw[vecteur] (0,1)--(3,1) node[nom,above]{$\vec{u}$};
  \draw[cherche] (3.6,2)--(5.4,-0.4) node[nom,right]{$\vec{v}$};
\end{tikzpicture}}
\rem{La direction est celle de la droite qui porte la flèche ; le sens dit de quel côté
est la pointe. Une même direction a deux sens : ne pas confondre les deux mots.}
\end{exo}

\begin{exo}[Vrai ou faux]
\cotefig[0.28]{%
\explique{On fait une figure : $EFGH$, puis $HGKL$ construit sur le côté
$[HG]$.}{le côté commun est $[HG]$}
\explique{Dans un parallélogramme, deux côtés opposés sont parallèles et de même
longueur.}{propriété du collège}}
{\begin{tikzpicture}[scale=0.5]
  \draw[fig] (0,0)--(3,0)--(4,2)--(1,2)--cycle;
  \draw[fig] (1,2)--(4,2)--(5,4)--(2,4)--cycle;
  \pt{0,0}{below left}{E} \pt{3,0}{below right}{F} \pt{4,2}{right}{G}
  \pt{1,2}{left}{H} \pt{5,4}{above right}{K} \pt{2,4}{above left}{L}
\end{tikzpicture}}
\cas{1)}
\explique{L'image de $H$ par la translation de vecteur $\vect{GF}$ est le point $H'$ tel
que $\vect{HH'}=\vect{GF}$.}{définition de l'image d'un point}
\explique{Dans $EFGH$, les côtés $[GF]$ et $[HE]$ sont parallèles et de même longueur ; de
$G$ vers $F$ et de $H$ vers $E$, on descend dans le même sens. Donc
$\vect{HE}=\vect{GF}$.}{même direction, même norme, même sens}
\explique{Donc $H'=E$ : l'affirmation est \res{vraie}.}{}
\cas{2)}
\explique{Dans $HGKL$, les côtés $[GK]$ et $[HL]$ sont parallèles.}{côtés opposés}
\explique{Donc $\vect{KG}$ et $\vect{HL}$ ont la même direction : l'affirmation est
\res{vraie}.}{la direction est celle de la droite}
\cas{3)}
\explique{$\vect{KG}$ descend de $K$ vers $G$ ; $\vect{HL}$ monte de $H$ vers $L$.}{on
regarde le sens}
\explique{Les sens sont contraires : $\vect{KG}=-\vect{HL}$. L'affirmation est
\res{fausse}.}{même direction, même norme, sens contraires : opposés}
\cas{4)}
\explique{Dans $EFGH$ : $[EF]$ et $[HG]$ sont parallèles, de même longueur, parcourus dans
le même sens, donc $\vect{EF}=\vect{HG}$.}{côtés opposés de $EFGH$}
\explique{Dans $HGKL$ : de même, $\vect{HG}=\vect{LK}$.}{côtés opposés de $HGKL$}
\explique{Donc $\vect{EF}=\vect{LK}$ : l'affirmation est \res{vraie}.}{deux vecteurs égaux
à un même troisième sont égaux}
\rem{Erreur fréquente (3) : conclure à l'égalité dès que les flèches sont parallèles et de
même longueur. Le sens compte : on vérifie où est la pointe.}
\end{exo}

\begin{exo}[Dans un parallélogramme]
\cotefig[0.28]{%
\explique{Dans le parallélogramme $RSTU$, les côtés opposés $[RS]$ et $[UT]$, $[RU]$ et
$[ST]$ sont parallèles et de même longueur.}{on fait le tour $R$, $S$, $T$, $U$}
\explique{Ses diagonales $[RT]$ et $[SU]$ se coupent en leur milieu $O$.}{propriété du
collège}}
{\begin{tikzpicture}[scale=0.5]
  \draw[fig] (0,0)--(4,0)--(5,2)--(1,2)--cycle;
  \draw[aide] (0,0)--(5,2) (4,0)--(1,2);
  \pt{0,0}{below left}{R} \pt{4,0}{below right}{S} \pt{5,2}{above right}{T}
  \pt{1,2}{above left}{U} \pt{2.5,1}{below}{O}
\end{tikzpicture}}
\cas{1)}
\explique{$\vect{RS}=\res{\vect{UT}}$.}{côté opposé, parcouru dans le même sens}
\explique{$\vect{TU}=\res{\vect{SR}}$.}{de $T$ vers $U$ : on va vers la gauche, comme de $S$
vers $R$}
\explique{$\vect{RU}=\res{\vect{ST}}$.}{les deux autres côtés opposés}
\cas{2)}
\explique{$O$ est le milieu de $[RT]$ : de $R$ à $O$, puis de $O$ à $T$, on fait le même
déplacement. Donc $\vect{RO}=\vect{O\res{T}}$.}{on continue dans la même direction}
\explique{De même, $O$ est le milieu de $[US]$ : $\vect{UO}=\vect{O\res{S}}$.}{}
\explique{Et $\vect{SO}=\vect{O\res{U}}$.}{}
\rem{Erreur fréquente : $\vect{RO}=\vect{OR}$. On revient alors en arrière : $\vect{OR}$
est l'opposé de $\vect{RO}$.}
\end{exo}

\partie{Translation}

\begin{exo}[Représentants d'origine donnée]
\cotefig[0.4]{%
\explique{On trace un triangle $DEF$ et un point $P$ hors du triangle.}{figure au choix}
\explique{Pour $\vect{DE}$ : depuis $P$, on trace une flèche parallèle à $(DE)$, de même
longueur que $DE$, pointée dans le même sens.}{même direction, même norme, même sens}
\explique{De même pour $\vect{EF}$ et pour $\vect{FD}$, toujours depuis $P$.}{les trois
flèches partent de $P$}
\explique{Sur un quadrillage, on compte les carreaux : $\vect{DE}$ fait $4$ carreaux à
droite, $\vect{EF}$ $3$ à gauche et $3$ vers le haut, $\vect{FD}$ $1$ à gauche et $3$
vers le bas.}{pour vérifier}}
{\begin{tikzpicture}[scale=0.42]
  \quadrillage{-0.4}{-1.4}{10.4}{5.4}
  \triangleDEF
  \pt{6,2}{right}{P}
  \draw[vecteur] (6,2)--(10,2) node[nom,above]{$\vect{DE}$};
  \draw[vecteurB] (6,2)--(3,5) node[nom,above right]{$\vect{EF}$};
  \draw[vecteurC] (6,2)--(5,-1) node[nom,left]{$\vect{FD}$};
\end{tikzpicture}}
\rem{Sans quadrillage : on trace la parallèle au côté passant par $P$ et on y reporte sa
longueur au compas, du bon côté.}
\end{exo}

\begin{exo}[Quatre points à placer]
\explique{On trace un triangle $ABC$ sur le quadrillage et on lit les déplacements :
$\vect{AC}$, $1$ carreau à droite et $3$ vers le haut ; $\vect{AB}$, $4$ à droite ;
$\vect{CB}$, $3$ à droite et $3$ vers le bas.}{méthode 2 : on reporte le déplacement}
\cas{Le point $M$ : $\vect{BM}=\vect{AC}$}
\cotefig{%
\explique{$B$ est l'origine : depuis $B$, on avance de $1$ à droite et $3$ vers le
haut.}{on reporte $\vect{AC}$ depuis $B$}
\explique{On arrive en $M$.}{}}
{\begin{tikzpicture}[scale=0.38]
  \quadrillage{-0.4}{-0.4}{5.4}{3.4}
  \triangleABC
  \draw[vecteur] (0,0)--(1,3) node[nom,left]{$\vect{AC}$};
  \draw[vecteur] (4,0)--(5,3) node[nom,right]{$\vect{AC}$};
  \pt{5,3}{above}{M}
\end{tikzpicture}}
\cas{Le point $N$ : $\vect{NB}=\vect{AC}$}
\cotefig{%
\explique{Cette fois, $B$ est l'\emph{extrémité} : la flèche doit arriver en $B$.}{$N$ est
l'origine}
\explique{Depuis $B$, on recule : $1$ à gauche et $3$ vers le bas. On arrive en $N$.}{on
fait le déplacement à l'envers}}
{\begin{tikzpicture}[scale=0.38]
  \quadrillage{-0.4}{-3.4}{5.4}{3.4}
  \triangleABC
  \draw[vecteur] (0,0)--(1,3) node[nom,left]{$\vect{AC}$};
  \draw[vecteur] (3,-3)--(4,0) node[nom,right]{$\vect{AC}$};
  \pt{3,-3}{below}{N}
\end{tikzpicture}}
\cas{Le point $P$ : $\vect{PC}=\vect{AB}$}
\cotefig{%
\explique{$C$ est l'extrémité : depuis $C$, on recule de $4$ carreaux vers la gauche. On
arrive en $P$.}{$P$ est l'origine}}
{\begin{tikzpicture}[scale=0.38]
  \quadrillage{-3.4}{-0.4}{4.4}{3.4}
  \triangleABC
  \draw[vecteurB] (0,0)--(4,0) node[nom,below]{$\vect{AB}$};
  \draw[vecteurB] (-3,3)--(1,3) node[nom,above]{$\vect{AB}$};
  \pt{-3,3}{left}{P}
\end{tikzpicture}}
\cas{Le point $Q$ : $\vect{AQ}=\vect{CB}$}
\cotefig{%
\explique{$A$ est l'origine : depuis $A$, on avance de $3$ à droite et $3$ vers le bas.
On arrive en $Q$.}{on reporte $\vect{CB}$ depuis $A$}}
{\begin{tikzpicture}[scale=0.38]
  \quadrillage{-0.4}{-3.4}{4.4}{3.4}
  \triangleABC
  \draw[vecteurC] (1,3)--(4,0) node[nom,above right]{$\vect{CB}$};
  \draw[vecteurC] (0,0)--(3,-3) node[nom,below left]{$\vect{CB}$};
  \pt{3,-3}{below}{Q}
\end{tikzpicture}}
\rem{Les points $Q$ et $N$ sont confondus : tous deux sont à $3$ carreaux à droite et $3$
vers le bas de $A$. Erreur fréquente : pour $N$ et $P$, partir du point connu en
avançant, alors qu'il est l'extrémité. On regarde toujours où est le point connu dans
l'égalité.}
\end{exo}

\begin{exo}[Origine ou extrémité]
\cotefig[0.36]{%
\explique{$\vec{u}$ fait $2$ carreaux à droite et $1$ vers le haut.}{on lit le
déplacement}
\explique{Origine $E$ : depuis $E$, on avance de $2$ à droite et $1$ vers le haut, et on
trace la flèche.}{la flèche part de $E$}
\explique{Extrémité $F$ : depuis $F$, on recule de $2$ à gauche et $1$ vers le bas ; puis on
trace la flèche de ce point vers $F$.}{la pointe est en $F$}
\explique{Extrémité $G$ : même geste depuis $G$.}{la pointe est en $G$}}
{\begin{tikzpicture}[scale=0.45]
  \quadrillage{-0.4}{-1.4}{7.4}{4.4}
  \draw[vecteur] (0,3)--(2,4) node[nom,above left]{$\vec{u}$};
  \pt{0,0}{below left}{E} \draw[vecteur] (0,0)--(2,1) node[nom,above left]{$\vec{u}$};
  \pt{6,3}{above right}{F} \draw[vecteur] (4,2)--(6,3) node[nom,above left]{$\vec{u}$};
  \pt{7,0}{right}{G} \draw[vecteur] (5,-1)--(7,0) node[nom,above left]{$\vec{u}$};
\end{tikzpicture}}
\rem{Les trois flèches sont trois représentants du même vecteur $\vec{u}$ : même
direction, même sens, même longueur. Seule leur position change.}
\end{exo}

\begin{exo}[Image d'un triangle]
\cotefig[0.36]{%
\explique{On construit le triangle à la règle et au compas : $[AB]$ de $5$~cm, puis un arc
de rayon $3$~cm autour de $A$ et un arc de rayon $4$~cm autour de $B$ ; ils se coupent
en $C$.}{figure ci-contre, réduite}
\explique{Depuis $A$, on reporte $\vect{BC}$ : on obtient $A'$.}{$\vect{AA'}=\vect{BC}$}
\explique{Depuis $B$, on reporte $\vect{BC}$ : on arrive en $C$, donc $B'=C$.}{$\vect{BB'}=\vect{BC}$}
\explique{Depuis $C$, on reporte $\vect{BC}$ : on obtient $C'$.}{$\vect{CC'}=\vect{BC}$}
\explique{On relie $A'$, $B'$ et $C'$.}{dans le même ordre que $A$, $B$, $C$}}
{\begin{tikzpicture}[scale=0.5]
  \quadrillage{-4.4}{-0.4}{4.4}{3.4}
  \draw[fig] (0,3)--(4,0)--(0,0)--cycle;
  \draw[figC] (-4,3)--(0,0)--(-4,0)--cycle;
  \draw[vecteurB] (4,0)--(0,0) node[nom,below]{$\vect{BC}$};
  \draw[vecteurB] (0,3)--(-4,3) node[nom,above]{$\vect{BC}$};
  \draw[vecteurB] (0,0)--(-4,0) node[nom,below]{$\vect{BC}$};
  \pt{0,3}{above right}{A} \pt{4,0}{below right}{B} \pt{0,0}{above right}{C}
  \pt{-4,3}{above left}{A'} \pt{-4,0}{below left}{C'}
  \node[font=\small,below] at (0,-0.2) {$B'$};
\end{tikzpicture}}
\explique{Les triangles $ABC$ et $A'B'C'$ sont \res{superposables} : $A'B'=5$~cm,
$B'C'=4$~cm et $A'C'=3$~cm.}{une translation conserve les longueurs (propriété 1)}
\rem{Le triangle glisse sans tourner : chaque côté de $A'B'C'$ est parallèle au côté
correspondant de $ABC$.}
\end{exo}

\begin{exo}[Neuf points]
\explique{Les points sont espacés de $2$ carreaux : $A$, $B$, $C$ en haut, $D$, $E$, $F$ au
milieu, $G$, $H$, $I$ en bas.}{on lit tout en carreaux}
\cas{1)}
\explique{$\vect{AB}$ : $2$ carreaux à droite. On cherche toutes les flèches qui joignent
deux points voisins sur une même ligne, de gauche à droite.}{même déplacement}
\explique{Les vecteurs égaux à $\vect{AB}$ sont $\res{\vect{BC}}$, $\res{\vect{DE}}$,
$\res{\vect{EF}}$, $\res{\vect{GH}}$ et $\res{\vect{HI}}$.}{cinq vecteurs}
\cas{2)}
\explique{$\vect{AE}$ : $2$ carreaux à droite et $2$ vers le bas. Depuis chaque point, on
regarde si ce déplacement arrive sur un autre point.}{$C$, $F$, $G$, $H$, $I$ sortent de
la figure}
\explique{$\vect{AE}=\res{\vect{BF}}$, $\vect{AE}=\res{\vect{DH}}$ et
$\vect{AE}=\res{\vect{EI}}$.}{trois vecteurs}
\cas{3)}
\explique{Un vecteur opposé à $\vect{AB}$ fait $2$ carreaux à gauche : par exemple
$\res{\vect{BA}}$ et $\res{\vect{CB}}$.}{on lit de droite à gauche}
\rem{Erreur fréquente : oublier $\vect{EF}$ ou $\vect{HI}$ en 1). On procède avec ordre,
ligne par ligne, pour n'en oublier aucun.}
\end{exo}

\partie{Parallélogramme et égalité de vecteurs}

\begin{exo}[Deux parallélogrammes]
\cotefig[0.3]{%
\explique{On trace $[EF]$, puis $EFGH$ d'un côté et $EFIJ$ de l'autre.}{$E$ et $F$ sont
communs}}
{\begin{tikzpicture}[scale=0.5]
  \draw[fig] (0,0)--(3,0)--(4,2)--(1,2)--cycle;
  \draw[fig] (0,0)--(3,0)--(2,-2)--(-1,-2)--cycle;
  \draw[aide] (4,2)--(2,-2) (1,2)--(-1,-2);
  \draw[vecteur] (0,0)--(3,0) node[nom,above]{$\vect{EF}$};
  \draw[vecteur] (1,2)--(4,2) node[nom,above]{$\vect{HG}$};
  \draw[vecteur] (-1,-2)--(2,-2) node[nom,below]{$\vect{JI}$};
  \pt{0,0}{left}{E} \pt{3,0}{right}{F} \pt{4,2}{above right}{G} \pt{1,2}{above left}{H}
  \pt{2,-2}{below right}{I} \pt{-1,-2}{below left}{J}
\end{tikzpicture}}
\cas{1)}
\explique{$EFGH$ est un parallélogramme, donc $\vect{EF}=\res{\vect{HG}}$.}{théorème 1 :
$ABCD$ parallélogramme donne $\vect{AB}=\vect{DC}$}
\explique{$EFIJ$ est un parallélogramme, donc $\vect{EF}=\res{\vect{JI}}$.}{même théorème,
avec $E$, $F$, $I$, $J$}
\cas{2)}
\explique{$\vect{HG}=\vect{EF}$ et $\vect{EF}=\vect{JI}$, donc $\vect{HG}=\vect{JI}$.}{on
enchaîne les égalités de 1)}
\explique{D'après le théorème 1, \res{$HGIJ$ est un parallélogramme}.}{$\vect{AB}=\vect{DC}$
donne $ABCD$ : ici $A=H$, $B=G$, $C=I$, $D=J$}
\rem{Pour nommer le parallélogramme, on lit les lettres dans l'ordre :
$\vect{HG}=\vect{JI}$ donne $H$, $G$, puis $I$, $J$ (la seconde flèche lue de la pointe
vers l'origine). On vérifie sur la figure qu'on fait bien le tour.}
\end{exo}

\begin{exo}[Autour d'un triangle]
\explique{On trace un triangle $DEF$ sur le quadrillage : $\vect{DE}$ fait $4$ carreaux à
droite ; $\vect{FE}$, $3$ à droite et $3$ vers le bas.}{pour compter les carreaux}
\cas{1)}
\cotefig{%
\explique{Depuis $F$, on reporte $\vect{DE}$ : $4$ carreaux à droite. C'est
$M$.}{$F$ est l'origine de $\vect{FM}$}
\explique{On obtient le parallélogramme \res{$DEMF$} : on joint $D$, $E$, $M$, $F$ dans cet
ordre.}{le tour de la figure}
\explique{On a construit $M$ tel que $\vect{DE}=\vect{FM}$.}{la construction}
\explique{D'après le théorème 1, $DEMF$ est un parallélogramme.}{$\vect{AB}=\vect{DC}$
donne $ABCD$}}
{\begin{tikzpicture}[scale=0.42]
  \quadrillage{-0.4}{-0.4}{5.4}{3.4}
  \triangleDEF
  \draw[figC] (4,0)--(5,3);
  \draw[vecteur] (0,0)--(4,0) node[nom,above]{$\vect{DE}$};
  \draw[vecteur] (1,3)--(5,3) node[nom,above]{$\vect{DE}$};
  \pt{5,3}{above right}{M}
\end{tikzpicture}}
\cas{2)}
\cotefig{%
\explique{Ici $F$ est l'extrémité : depuis $F$, on recule de $4$ carreaux vers la gauche.
C'est $N$.}{$\vect{NF}$ arrive en $F$}
\explique{On obtient le parallélogramme \res{$DEFN$}.}{le tour de la figure}
\explique{On a construit $N$ tel que $\vect{DE}=\vect{NF}$.}{la construction}
\explique{D'après le théorème 1, $DEFN$ est un parallélogramme.}{$\vect{AB}=\vect{DC}$
donne $ABCD$}}
{\begin{tikzpicture}[scale=0.42]
  \quadrillage{-3.4}{-0.4}{4.4}{3.4}
  \triangleDEF
  \draw[figC] (0,0)--(-3,3);
  \draw[vecteur] (0,0)--(4,0) node[nom,below]{$\vect{DE}$};
  \draw[vecteur] (-3,3)--(1,3) node[nom,above]{$\vect{DE}$};
  \pt{-3,3}{above left}{N}
\end{tikzpicture}}
\cas{3)}
\cotefig{%
\explique{L'image $K$ de $D$ vérifie $\vect{DK}=\vect{FE}$ : depuis $D$, on avance de $3$ à
droite et $3$ vers le bas.}{définition de l'image}
\explique{On obtient le parallélogramme \res{$FEKD$}.}{le tour de la figure}
\explique{On a construit $K$ tel que $\vect{FE}=\vect{DK}$.}{la construction}
\explique{D'après le théorème 1, $FEKD$ est un parallélogramme.}{$\vect{AB}=\vect{DC}$
donne $ABCD$, avec $A=F$, $B=E$}}
{\begin{tikzpicture}[scale=0.42]
  \quadrillage{-0.4}{-3.4}{4.4}{3.4}
  \triangleDEF
  \draw[figC] (4,0)--(3,-3);
  \draw[vecteurB] (1,3)--(4,0) node[nom,above right]{$\vect{FE}$};
  \draw[vecteurB] (0,0)--(3,-3) node[nom,below left]{$\vect{FE}$};
  \pt{3,-3}{below}{K}
\end{tikzpicture}}
\cas{4)}
\cotefig{%
\explique{On cherche le point $X$ tel que $\vect{XE}=\vect{FD}$ : depuis $E$, on recule du
déplacement de $\vect{FD}$ ; on arrive en $M$.}{$E$ est l'extrémité}
\explique{Justification : $DEMF$ est un parallélogramme ; en faisant le tour depuis $M$
dans l'autre sens, c'est $MEDF$.}{le même quadrilatère, lu autrement}
\explique{D'après le théorème 1, $\vect{ME}=\vect{FD}$ : le point cherché est
\res{$M$}.}{$\vect{AB}=\vect{DC}$ avec $A=M$, $B=E$, $C=D$, $D=F$}}
{\begin{tikzpicture}[scale=0.42]
  \quadrillage{-0.4}{-0.4}{5.4}{3.4}
  \triangleDEF
  \draw[figC] (4,0)--(5,3)--(1,3);
  \draw[vecteurC] (1,3)--(0,0) node[nom,left]{$\vect{FD}$};
  \draw[vecteurC] (5,3)--(4,0) node[nom,right]{$\vect{FD}$};
  \pt{5,3}{above right}{M}
\end{tikzpicture}}
\rem{Erreur fréquente en 2) : reporter $\vect{DE}$ en partant de $F$. Dans $\vect{NF}$, $F$
est la pointe de la flèche : on part de $F$ en reculant.}
\end{exo}

\begin{exo}[Un parallélogramme de plus]
\cas{1) et 2)}
\cotefig{%
\explique{On trace le parallélogramme $RSTU$.}{1) : figure au choix}
\explique{$RSUV$ parallélogramme équivaut à $\vect{RS}=\vect{VU}$.}{théorème 1 : il
faut $\vect{AB}=\vect{DC}$}
\explique{$U$ est l'extrémité : depuis $U$, on recule du déplacement de $\vect{RS}$, ici $3$
carreaux vers la gauche. C'est $V$.}{on reporte $\vect{RS}$ pour qu'il arrive en $U$}
\explique{On trace $RSUV$ en joignant $R$, $S$, $U$, $V$.}{le tour de la figure}}
{\begin{tikzpicture}[scale=0.45]
  \quadrillage{-2.4}{-0.4}{4.4}{2.4}
  \draw[fig] (0,0)--(3,0)--(4,2)--(1,2)--cycle;
  \draw[figC] (3,0)--(1,2) (0,0)--(-2,2);
  \draw[vecteur] (0,0)--(3,0) node[nom,below]{$\vect{RS}$};
  \draw[vecteur] (-2,2)--(1,2) node[nom,above]{$\vect{RS}$};
  \pt{0,0}{below left}{R} \pt{3,0}{below right}{S} \pt{4,2}{above right}{T}
  \pt{1,2}{above}{U} \pt{-2,2}{above left}{V}
\end{tikzpicture}}
\cas{3)}
\explique{$RSTU$ est un parallélogramme, donc $\vect{RS}=\res{\vect{UT}}$.}{théorème 1}
\explique{$RSUV$ est un parallélogramme, donc $\vect{RS}=\res{\vect{VU}}$.}{théorème 1}
\rem{$U$ est le milieu de $[VT]$ : $\vect{VU}=\vect{UT}$. On le démontrera dans la partie V
(exercice 21).}
\end{exo}

\begin{exo}[Vrai ou faux]
\cas{1)}
\explique{$F$ est l'image de $C$ par la translation de vecteur $\vect{DE}$, donc
$\vect{CF}=\vect{DE}$.}{le vecteur va du point à son image}
\explique{$\vect{FC}$ est l'opposé de $\vect{CF}$ : $\vect{FC}=-\vect{DE}$, et non
$\vect{DE}$.}{on a échangé les lettres}
\explique{L'affirmation est \res{fausse} (sauf si $D=E$).}{}
\cas{2)}
\explique{$\vect{KL}=\vect{NM}$ est de la forme $\vect{AB}=\vect{DC}$ avec $A=K$, $B=L$,
$C=M$, $D=N$.}{on identifie les lettres}
\explique{D'après le théorème 1, $KLMN$ est un parallélogramme : l'affirmation est
\res{vraie}.}{éventuellement aplati, si les points sont alignés}
\cas{3)}
\explique{$ABCD$ est un parallélogramme ; lu dans l'autre sens, c'est $ADCB$, donc
$\vect{AD}=\vect{BC}$.}{théorème 1 avec $A$, $D$, $C$, $B$}
\explique{L'image de $B$ est le point $B'$ tel que $\vect{BB'}=\vect{AD}$, soit
$\vect{BB'}=\vect{BC}$. Donc $B'=C$ : l'affirmation est \res{vraie}.}{définition de
l'image}
\rem{Une figure à main levée aide à décider ; pour justifier, on cite le théorème 1 ou la
définition de la translation.}
\end{exo}

\partie{Somme de deux vecteurs}

\begin{exo}[Origine imposée]
\explique{Pour tracer $\vec{w}=\vec{u}+\vec{v}$ depuis $T$, on trace un représentant de
$\vec{u}$ d'origine $T$, puis un représentant de $\vec{v}$ à la suite.}{méthode 4 : l'un à
la suite de l'autre}
\cas{Figure de gauche}
\cotefig{%
\explique{$\vec{u}$ : $1$ carreau à droite et $2$ vers le haut ; on le trace depuis
$T$.}{on lit, puis on reporte}
\explique{$\vec{v}$ : $2$ carreaux à droite et $1$ vers le bas ; on le trace à la suite de
$\vec{u}$.}{origine de $\vec{v}$ sur la pointe de $\vec{u}$}
\explique{$\vec{w}$ va du départ de $\vec{u}$ (en $T$) à l'arrivée de $\vec{v}$ : $3$
carreaux à droite et $1$ vers le haut.}{départ, arrivée}}
{\begin{tikzpicture}[scale=0.5]
  \quadrillage{-0.4}{-0.4}{5.4}{4.4}
  \draw[vecteur] (0,2)--(1,4) node[nom,left]{$\vec{u}$};
  \draw[vecteurB] (3,4)--(5,3) node[nom,above right]{$\vec{v}$};
  \draw[vecteur] (1,1)--(2,3) node[nom,left]{$\vec{u}$};
  \draw[vecteurB] (2,3)--(4,2) node[nom,above right]{$\vec{v}$};
  \draw[cherche] (1,1)--(4,2) node[nom,below right]{$\vec{w}$};
  \pt{1,1}{below left}{T}
\end{tikzpicture}}
\cas{Figure de droite}
\cotefig{%
\explique{$\vec{u}$ : $2$ carreaux à droite et $1$ vers le bas ; on le trace depuis
$T$.}{on lit, puis on reporte}
\explique{$\vec{v}$ : $1$ carreau à gauche et $2$ vers le bas ; on le trace à la suite de
$\vec{u}$.}{origine de $\vec{v}$ sur la pointe de $\vec{u}$}
\explique{$\vec{w}$ va de $T$ à l'arrivée de $\vec{v}$ : $1$ carreau à droite et $3$ vers le
bas.}{départ, arrivée}}
{\begin{tikzpicture}[scale=0.5]
  \quadrillage{-0.4}{-0.4}{5.4}{4.4}
  \draw[vecteur] (0,4)--(2,3) node[nom,below left]{$\vec{u}$};
  \draw[vecteurB] (2,2)--(1,0) node[nom,left]{$\vec{v}$};
  \draw[vecteur] (3,4)--(5,3) node[nom,above right]{$\vec{u}$};
  \draw[vecteurB] (5,3)--(4,1) node[nom,right]{$\vec{v}$};
  \draw[cherche] (3,4)--(4,1) node[nom,left]{$\vec{w}$};
  \pt{3,4}{above}{T}
\end{tikzpicture}}
\rem{On peut vérifier par le calcul des carreaux : à gauche, $1+2=3$ à droite et $2-1=1$
vers le haut. Erreur fréquente : tracer $\vec{u}$ et $\vec{v}$ depuis $T$ tous les deux ; il
faut les mettre bout à bout.}
\end{exo}

\begin{exo}[Trois points à construire]
\explique{On trace un triangle $ABC$ sur le quadrillage. Pour chaque point, on met les deux
vecteurs bout à bout depuis l'origine imposée.}{méthode 4}
\cas{Le point $D$ : $\vect{BD}=\vect{BC}+\vect{CA}$}
\cotefig{%
\explique{$\vect{BC}$ et $\vect{CA}$ sont déjà bout à bout : de $B$ on va en $C$, puis de $C$
en $A$.}{la pointe du premier est l'origine du second}
\explique{La somme va du départ, $B$, à l'arrivée, $A$ : $\vect{BD}=\vect{BA}$, donc
\res{$D=A$}.}{départ, arrivée}}
{\begin{tikzpicture}[scale=0.42]
  \quadrillage{-0.4}{-0.4}{4.4}{3.4}
  \draw[fig] (0,0)--(4,0)--(1,3)--cycle;
  \draw[vecteurB] (4,0)--(1,3) node[nom,above right]{$\vect{BC}$};
  \draw[vecteurC] (1,3)--(0,0) node[nom,left]{$\vect{CA}$};
  \draw[cherche] (4,0)--(0,0) node[nom,below]{$\vect{BD}$};
  \pt{0,0}{left}{A} \pt{4,0}{right}{B} \pt{1,3}{above}{C}
  \node[font=\small,below left] at (0,-0.4) {$D$};
\end{tikzpicture}}
\cas{Le point $E$ : $\vect{BE}=\vect{BA}+\vect{BC}$}
\cotefig{%
\explique{On trace $\vect{BA}$, puis un représentant de $\vect{BC}$ à la suite, depuis $A$
: $3$ carreaux à gauche et $3$ vers le haut.}{l'un à la suite de l'autre}
\explique{$\vect{BE}$ va de $B$ à l'arrivée de ce représentant : c'est le point
$E$.}{départ, arrivée}
\explique{$BAEC$ est un parallélogramme : $\vect{BC}=\vect{AE}$.}{les deux vecteurs ont la
même origine $B$ : la somme suit la diagonale}}
{\begin{tikzpicture}[scale=0.42]
  \quadrillage{-3.4}{-0.4}{4.4}{3.4}
  \draw[fig] (0,0)--(4,0)--(1,3)--cycle;
  \draw[aide] (-3,3)--(1,3);
  \draw[vecteur] (4,0)--(0,0) node[nom,below]{$\vect{BA}$};
  \draw[vecteurB] (0,0)--(-3,3) node[nom,below left]{$\vect{BC}$};
  \draw[cherche] (4,0)--(-3,3) node[nom,above,pos=0.35]{$\vect{BE}$};
  \pt{0,0}{below}{A} \pt{4,0}{below right}{B} \pt{1,3}{above right}{C}
  \pt{-3,3}{above}{E}
\end{tikzpicture}}
\cas{Le point $F$ : $\vect{CF}=\vect{CB}+\vect{AB}$}
\cotefig{%
\explique{On trace $\vect{CB}$, puis un représentant de $\vect{AB}$ à la suite, depuis $B$
: $4$ carreaux à droite.}{l'un à la suite de l'autre}
\explique{$\vect{CF}$ va de $C$ à l'arrivée de ce représentant : c'est le point
$F$.}{départ, arrivée}}
{\begin{tikzpicture}[scale=0.42]
  \quadrillage{-0.4}{-0.4}{8.4}{3.4}
  \draw[fig] (0,0)--(4,0)--(1,3)--cycle;
  \draw[vecteurC] (1,3)--(4,0) node[nom,left]{$\vect{CB}$};
  \draw[vecteur] (4,0)--(8,0) node[nom,below]{$\vect{AB}$};
  \draw[cherche] (1,3)--(8,0) node[nom,above right]{$\vect{CF}$};
  \pt{0,0}{below left}{A} \pt{4,0}{below}{B} \pt{1,3}{above}{C} \pt{8,0}{below}{F}
\end{tikzpicture}}
\rem{Pour $D$, aucun tracé n'est nécessaire : quand la seconde lettre du premier vecteur
est la première du second, la somme se lit directement.}
\end{exo}

\begin{exo}[Somme et différence sur un quadrillage]
\explique{$\vect{AM}=\vec{u}+\vec{v}$ : on trace $\vec{v}$ à la suite de $\vec{u}$.
$\vect{AN}=\vec{u}-\vec{v}=\vec{u}+(-\vec{v})$ : on trace $-\vec{v}$ à la suite de
$\vec{u}$.}{méthode 4 : la différence est la somme avec l'opposé}
\cas{Figure de gauche}
\cotefig[0.46]{%
\explique{$\vec{u}$ : $3$ à droite et $1$ vers le haut ; $\vec{v}$ : $1$ à droite et $2$
vers le haut.}{on lit les déplacements}
\explique{Point $M$ : depuis la pointe de $\vec{u}$, on trace $\vec{v}$ ; $M$ est à
l'arrivée.}{$\vec{u}+\vec{v}$ : $4$ à droite, $3$ vers le haut}
\explique{Point $N$ : depuis la pointe de $\vec{u}$, on trace $-\vec{v}$ ($1$ à gauche et $2$
vers le bas) ; $N$ est à l'arrivée.}{$\vec{u}-\vec{v}$ : $2$ à droite, $1$ vers le bas}}
{\begin{tikzpicture}[scale=0.42]
  \quadrillage{-0.4}{-0.4}{5.4}{4.4}
  \draw[vecteur] (1,1)--(4,2) node[nom,below right]{$\vec{u}$};
  \draw[vecteurB] (1,1)--(2,3) node[nom,left]{$\vec{v}$};
  \draw[vecteurB] (4,2)--(5,4) node[nom,right]{$\vec{v}$};
  \draw[cherche] (1,1)--(5,4) node[nom,above left]{$\vect{AM}$};
  \pt{1,1}{below left}{A} \pt{5,4}{above}{M}
\end{tikzpicture}\hspace{4pt}%
\begin{tikzpicture}[scale=0.42]
  \quadrillage{-0.4}{-0.4}{5.4}{4.4}
  \draw[vecteur] (1,1)--(4,2) node[nom,above left]{$\vec{u}$};
  \draw[vecteurB] (1,1)--(2,3) node[nom,left]{$\vec{v}$};
  \draw[vecteurB] (4,2)--(3,0) node[nom,right]{$-\vec{v}$};
  \draw[cherche] (1,1)--(3,0) node[nom,below left]{$\vect{AN}$};
  \pt{1,1}{left}{A} \pt{3,0}{below}{N}
\end{tikzpicture}}
\cas{Figure de droite}
\cotefig[0.46]{%
\explique{$\vec{u}$ : $3$ à droite ; $\vec{v}$ : $2$ vers le bas.}{on lit les
déplacements}
\explique{Point $M$ : depuis la pointe de $\vec{u}$, on trace $\vec{v}$, $2$ vers le bas ;
$M$ est à l'arrivée.}{$\vec{u}+\vec{v}$ : $3$ à droite, $2$ vers le bas}
\explique{Point $N$ : depuis la pointe de $\vec{u}$, on trace $-\vec{v}$, $2$ vers le haut ;
$N$ est à l'arrivée.}{$\vec{u}-\vec{v}$ : $3$ à droite, $2$ vers le haut}}
{\begin{tikzpicture}[scale=0.42]
  \quadrillage{-0.4}{-0.4}{5.4}{4.4}
  \draw[vecteur] (1,2)--(4,2) node[nom,above]{$\vec{u}$};
  \draw[vecteurB] (1,2)--(1,0) node[nom,left]{$\vec{v}$};
  \draw[vecteurB] (4,2)--(4,0) node[nom,right]{$\vec{v}$};
  \draw[cherche] (1,2)--(4,0) node[nom,below left]{$\vect{AM}$};
  \pt{1,2}{left}{A} \pt{4,0}{below}{M}
\end{tikzpicture}\hspace{4pt}%
\begin{tikzpicture}[scale=0.42]
  \quadrillage{-0.4}{-0.4}{5.4}{4.4}
  \draw[vecteur] (1,2)--(4,2) node[nom,below]{$\vec{u}$};
  \draw[vecteurB] (1,2)--(1,0) node[nom,left]{$\vec{v}$};
  \draw[vecteurB] (4,2)--(4,4) node[nom,right]{$-\vec{v}$};
  \draw[cherche] (1,2)--(4,4) node[nom,above left]{$\vect{AN}$};
  \pt{1,2}{left}{A} \pt{4,4}{above}{N}
\end{tikzpicture}}
\rem{Erreur fréquente : relier les pointes de $\vec{u}$ et de $\vec{v}$, tracés tous deux
depuis $A$. On reporte toujours le second vecteur à la suite du premier.}
\end{exo}

\begin{exo}[Quand c'est l'origine qui manque]
\explique{Le point connu $A$ est l'\emph{extrémité} : la flèche doit arriver en $A$. On
part donc de $A$ et on fait les déplacements à l'envers.}{$N$ est l'origine}
\cas{1)}
\cotefig{%
\explique{$\vec{u}$ : $1$ à droite et $2$ vers le haut ; $\vec{v}$ : $2$ à droite et $1$
vers le bas.}{on lit les déplacements}
\explique{Depuis $A$, on recule de $\vec{v}$ ($2$ à gauche et $1$ vers le haut), puis de
$\vec{u}$ ($1$ à gauche et $2$ vers le bas) : on arrive en $N$.}{on remonte la chaîne
depuis la fin}
\explique{Vérification : depuis $N$, $\vec{u}$ puis $\vec{v}$ arrivent bien en $A$.}{les
vecteurs bout à bout de $N$ à $A$}}
{\begin{tikzpicture}[scale=0.5]
  \quadrillage{-0.4}{-0.4}{5.4}{4.4}
  \draw[vecteur] (0,1)--(1,3) node[nom,left]{$\vec{u}$};
  \draw[vecteurB] (2,4)--(4,3) node[nom,above right]{$\vec{v}$};
  \draw[vecteur] (1,1)--(2,3) node[nom,left]{$\vec{u}$};
  \draw[vecteurB] (2,3)--(4,2) node[nom,above right]{$\vec{v}$};
  \draw[cherche] (1,1)--(4,2) node[nom,below right]{$\vect{NA}$};
  \pt{4,2}{right}{A} \pt{1,1}{below left}{N}
\end{tikzpicture}}
\cas{2)}
\cotefig{%
\explique{$\vec{u}$ : $2$ à droite et $1$ vers le haut ; $-\vec{v}$ : $1$ à gauche et $2$
vers le bas.}{$\vec{u}-\vec{v}=\vec{u}+(-\vec{v})$}
\explique{Depuis $A$, on recule de $-\vec{v}$ ($1$ à droite et $2$ vers le haut), puis de
$\vec{u}$ ($2$ à gauche et $1$ vers le bas) : on arrive en $N$.}{on remonte la chaîne}
\explique{Vérification : depuis $N$, $\vec{u}$ puis $-\vec{v}$ arrivent bien en $A$.}{}}
{\begin{tikzpicture}[scale=0.5]
  \quadrillage{-0.4}{-0.4}{5.4}{4.4}
  \draw[vecteur] (0,3)--(2,4) node[nom,above left]{$\vec{u}$};
  \draw[vecteurB] (4,0)--(5,2) node[nom,right]{$\vec{v}$};
  \draw[vecteur] (1,3)--(3,4) node[nom,above]{$\vec{u}$};
  \draw[vecteurB] (3,4)--(2,2) node[nom,right]{$-\vec{v}$};
  \draw[cherche] (1,3)--(2,2) node[nom,below left]{$\vect{NA}$};
  \pt{2,2}{below right}{A} \pt{1,3}{left}{N}
\end{tikzpicture}}
\rem{Autre façon : on construit d'abord le représentant de la somme (ou de la différence)
d'origine $A$, puis $N$ est le point obtenu en reculant de ce déplacement depuis $A$. À
gauche, $\vec{u}+\vec{v}$ fait $3$ à droite et $1$ vers le haut, donc $N$ est $3$ carreaux
à gauche et $1$ plus bas que $A$.}
\end{exo}

\begin{exo}[Avec des longueurs]
\explique{On trace $[AB]$ de $6$~cm ; $C$ est à l'intersection d'un arc de rayon $5$~cm
centré en $A$ et d'un arc de rayon $4$~cm centré en $B$.}{figures réduites de moitié}
\cas{1)}
\cotefig[0.4]{%
\explique{On trace un représentant de $\vect{AC}$ à la suite de $\vect{AB}$, depuis $B$ :
parallèle à $(AC)$, de longueur $5$~cm, dans le même sens. Il arrive en
$D$.}{l'un à la suite de l'autre}
\explique{$\vect{AD}$ va du départ de $\vect{AB}$ à l'arrivée de ce
représentant.}{départ, arrivée}}
{\begin{tikzpicture}[scale=0.42]
  \draw[aide] (0,0) ++(25:5) arc (25:60:5) (6,0) ++(105:4) arc (105:145:4);
  \draw[fig] (0,0)--(6,0)--(3.75,3.307)--cycle;
  \draw[vecteur] (0,0)--(6,0) node[nom,below]{$\vect{AB}$};
  \draw[vecteurB] (6,0)--(9.75,3.307) node[nom,right]{$\vect{AC}$};
  \draw[cherche] (0,0)--(9.75,3.307) node[nom,above,pos=0.6]{$\vect{AD}$};
  \pt{0,0}{below left}{A} \pt{6,0}{below right}{B} \pt{3.75,3.307}{above}{C}
  \pt{9.75,3.307}{above right}{D}
\end{tikzpicture}}
\cas{2)}
\cotefig[0.4]{%
\explique{$\vect{BC}-\vect{BA}=\vect{BC}+\vect{AB}$.}{l'opposé de $\vect{BA}$ est
$\vect{AB}$}
\explique{On trace $\vect{BC}$, puis un représentant de $\vect{AB}$ à la suite, depuis $C$ :
parallèle à $(AB)$, $6$~cm, vers la droite. Il arrive en $E$.}{l'un à la suite de
l'autre}
\explique{On constate que \res{$E$ est confondu avec $D$}.}{à expliquer}
\explique{En effet, $\vect{BC}+\vect{AB}=\vect{AB}+\vect{BC}$ : bout à bout, de $A$ à $C$.
Donc $\vect{BE}=\vect{AC}$, comme $\vect{BD}$.}{l'ordre ne compte pas}}
{\begin{tikzpicture}[scale=0.42]
  \draw[fig] (0,0)--(6,0)--(3.75,3.307)--cycle;
  \draw[vecteurB] (6,0)--(3.75,3.307) node[nom,left]{$\vect{BC}$};
  \draw[vecteur] (3.75,3.307)--(9.75,3.307) node[nom,above]{$\vect{AB}$};
  \draw[cherche] (6,0)--(9.75,3.307) node[nom,right]{$\vect{BE}$};
  \pt{0,0}{below left}{A} \pt{6,0}{below right}{B} \pt{3.75,3.307}{above left}{C}
  \pt{9.75,3.307}{above right}{E}
\end{tikzpicture}}
\cas{3)}
\explique{Par construction, $\vect{BD}=\vect{AC}$, soit $\vect{AC}=\vect{BD}$.}{la
construction de 1)}
\explique{D'après le théorème 1, $ACDB$ est un parallélogramme, c'est-à-dire
\res{$ABDC$}.}{$\vect{AB}=\vect{DC}$ donne $ABCD$, avec $A$, $C$, $D$, $B$}
\explique{Ses côtés mesurent $6$~cm et $5$~cm, et ce n'est pas un rectangle.}{$AB$ et $AC$ ne
sont pas perpendiculaires}
\rem{Deux vecteurs de même origine $A$ : leur somme suit la diagonale $[AD]$ du
parallélogramme construit sur eux.}
\end{exo}

\begin{exo}[Trois vecteurs]
\explique{$\vec{u}$ : $1$ à droite et $1$ vers le haut ; $\vec{v}$ : $2$ à droite et $1$
vers le bas ; $\vec{w}$ : $1$ à gauche et $2$ vers le haut.}{on lit les trois
déplacements}
\cas{1)}
\cotefig[0.42]{%
\explique{Depuis $A$, on trace $\vec{u}$ ; à sa suite, on trace $\vec{v}$. $\vec{r}$ va de
$A$ à l'arrivée de $\vec{v}$ : $3$ carreaux à droite.}{$\vec{r}=\vec{u}+\vec{v}$}
\explique{À la suite du même $\vec{u}$, on trace $\vec{w}$. $\vec{s}$ va de $A$ à l'arrivée
de $\vec{w}$ : $3$ carreaux vers le haut.}{$\vec{s}=\vec{u}+\vec{w}$}}
{\begin{tikzpicture}[scale=0.45]
  \quadrillage{-0.4}{-0.4}{8.4}{4.4}
  \draw[vecteur] (2,3)--(3,4) node[nom,above left]{$\vec{u}$};
  \draw[vecteurB] (4,4)--(6,3) node[nom,above right]{$\vec{v}$};
  \draw[vecteurC] (8,1)--(7,3) node[nom,right]{$\vec{w}$};
  \draw[vecteur] (1,1)--(2,2) node[nom,below right]{$\vec{u}$};
  \draw[vecteurB] (2,2)--(4,1) node[nom,above right]{$\vec{v}$};
  \draw[vecteurC] (2,2)--(1,4) node[nom,right]{$\vec{w}$};
  \draw[cherche] (1,1)--(4,1) node[nom,below]{$\vec{r}$};
  \draw[cherche] (1,1)--(1,4) node[nom,left]{$\vec{s}$};
  \pt{1,1}{below left}{A} \pt{5,0}{below}{B}
\end{tikzpicture}}
\cas{2)}
\cotefig[0.42]{%
\explique{Depuis $B$, on trace $\vec{v}$, puis $\vec{w}$ à sa suite. $\vec{t}$ va de $B$ à
l'arrivée de $\vec{w}$ : $1$ carreau à droite et $1$ vers le haut.}{$\vec{t}=\vec{v}+\vec{w}$}}
{\begin{tikzpicture}[scale=0.45]
  \quadrillage{-0.4}{-1.4}{8.4}{4.4}
  \draw[vecteur] (2,3)--(3,4) node[nom,above left]{$\vec{u}$};
  \draw[vecteurB] (4,4)--(6,3) node[nom,above right]{$\vec{v}$};
  \draw[vecteurC] (8,1)--(7,3) node[nom,right]{$\vec{w}$};
  \draw[vecteurB] (5,0)--(7,-1) node[nom,below left]{$\vec{v}$};
  \draw[vecteurC] (7,-1)--(6,1) node[nom,right]{$\vec{w}$};
  \draw[cherche] (5,0)--(6,1) node[nom,left]{$\vec{t}$};
  \pt{1,1}{below left}{A} \pt{5,0}{below left}{B}
\end{tikzpicture}}
\rem{Vérification par les carreaux : $\vec{t}$, $2-1=1$ à droite et $-1+2=1$ vers le haut.
On remarque que $\vec{t}=\vec{u}$.}
\end{exo}

\partie{Milieu d'un segment}

\begin{exo}[Un milieu à démontrer]
\cas{1) et 2)}
\cotefig[0.36]{%
\explique{On trace le parallélogramme $KLMN$.}{1) : figure au choix}
\explique{$KLPM$ parallélogramme équivaut à $\vect{KL}=\vect{MP}$ : depuis $M$, on reporte
$\vect{KL}$ ; c'est $P$.}{théorème 1 : $\vect{AB}=\vect{DC}$}
\explique{On trace $KLPM$ en joignant $K$, $L$, $P$, $M$.}{le tour de la figure}}
{\begin{tikzpicture}[scale=0.5]
  \quadrillage{-0.4}{-0.4}{7.4}{2.4}
  \draw[fig] (0,0)--(3,0)--(4,2)--(1,2)--cycle;
  \draw[figC] (3,0)--(7,2)--(4,2);
  \draw[vecteur] (0,0)--(3,0) node[nom,below]{$\vect{KL}$};
  \draw[vecteur] (4,2)--(7,2) node[nom,above]{$\vect{KL}$};
  \pt{0,0}{below left}{K} \pt{3,0}{below}{L} \pt{4,2}{above}{M}
  \pt{1,2}{above left}{N} \pt{7,2}{above right}{P}
\end{tikzpicture}}
\cas{3)}
\explique{$KLMN$ est un parallélogramme, donc $\vect{NM}=\vect{KL}$.}{côtés opposés
(théorème 1)}
\explique{$KLPM$ est un parallélogramme, donc $\vect{KL}=\vect{MP}$.}{côtés opposés
(théorème 1)}
\explique{Ainsi $\vect{NM}=\vect{MP}$ : \res{$M$ est le milieu de $[NP]$}.}{théorème 2 :
$\vect{XM}=\vect{MY}$}
\rem{Le vecteur $\vect{KL}$ sert de relais entre les deux parallélogrammes. On vise
$\vect{NM}=\vect{MP}$ : le point $M$ est l'extrémité du premier vecteur et l'origine du
second.}
\end{exo}

\begin{exo}[Une différence, un parallélogramme]
\cas{1)}
\cotefig{%
\explique{On trace un triangle $ABC$ sur le quadrillage.}{figure au choix}
\explique{$\vect{AC}-\vect{AB}=\vect{AC}+\vect{BA}$.}{la différence est la somme avec
l'opposé : $-\vect{AB}=\vect{BA}$}
\explique{On trace $\vect{AC}$, puis un représentant de $\vect{BA}$ à la suite, depuis $C$ :
$4$ carreaux à gauche. Il arrive en $D$.}{l'un à la suite de l'autre}
\explique{Autre lecture : $\vect{BA}+\vect{AC}$, bout à bout, va de $B$ à $C$. Donc
$\vect{AD}=\vect{BC}$.}{l'ordre ne compte pas}
\explique{D'après le théorème 1, \res{$ADCB$ est un parallélogramme}.}{$\vect{AB}=\vect{DC}$
avec $A$, $D$, $C$, $B$}
\explique{$[AC]$ et $[DB]$ sont ses diagonales : \res{elles ont le même milieu}.}{les
diagonales d'un parallélogramme se coupent en leur milieu}}
{\begin{tikzpicture}[scale=0.45]
  \quadrillage{-3.4}{-0.4}{4.4}{3.4}
  \draw[fig] (0,0)--(4,0)--(1,3)--cycle;
  \draw[figC] (0,0)--(-3,3);
  \draw[aide] (-3,3)--(4,0);
  \draw[vecteur] (0,0)--(1,3) node[nom,right]{$\vect{AC}$};
  \draw[vecteurB] (1,3)--(-3,3) node[nom,above]{$\vect{BA}$};
  \draw[cherche] (0,0)--(-3,3) node[nom,below left]{$\vect{AD}$};
  \pt{0,0}{below left}{A} \pt{4,0}{below right}{B} \pt{1,3}{above right}{C}
  \pt{-3,3}{above left}{D}
\end{tikzpicture}}
\cas{2)}
\cotefig{%
\explique{$\vect{BA}+\vect{CB}=\vect{CB}+\vect{BA}$.}{l'ordre ne compte pas}
\explique{$\vect{CB}$ et $\vect{BA}$ sont bout à bout : de $C$ à $B$, puis de $B$ à $A$.
Donc $\vect{CE}=\vect{CA}$.}{départ $C$, arrivée $A$}
\explique{On constate que \res{$E$ est confondu avec $A$}.}{}}
{\begin{tikzpicture}[scale=0.45]
  \quadrillage{-0.4}{-0.4}{4.4}{3.4}
  \draw[fig] (0,0)--(4,0)--(1,3)--cycle;
  \draw[vecteurC] (1,3)--(4,0) node[nom,above right]{$\vect{CB}$};
  \draw[vecteur] (4,0)--(0,0) node[nom,below]{$\vect{BA}$};
  \draw[cherche] (1,3)--(0,0) node[nom,left]{$\vect{CE}$};
  \pt{0,0}{above left}{A} \pt{4,0}{below right}{B} \pt{1,3}{above}{C}
  \node[font=\small,below left] at (0,-0.1) {$E$};
\end{tikzpicture}}
\rem{Erreur fréquente en 1) : tracer $\vect{AC}$ et $\vect{AB}$ depuis $A$ et relier leurs
pointes. Pour une différence, on met l'\emph{opposé} à la suite.}
\end{exo}

\begin{exo}[Trois points à construire]
\cas{1)}
\cotefig{%
\explique{On garde le même triangle $ABC$ ; chaque question donne un point $M$
différent.}{trois figures}
\explique{$\vect{AC}-\vect{AB}=\vect{AC}+\vect{BA}$, soit $\vect{BA}+\vect{AC}$ : bout à
bout, de $B$ à $C$.}{l'opposé de $\vect{AB}$ est $\vect{BA}$}
\explique{Donc $\vect{AM}=\vect{BC}$ : depuis $A$, on reporte $\vect{BC}$, $3$ à gauche et
$3$ vers le haut.}{$ABCM$ est un parallélogramme}
\explique{C'est le point $D$ de l'exercice 22.}{même égalité}}
{\begin{tikzpicture}[scale=0.4]
  \quadrillage{-3.4}{-0.4}{4.4}{3.4}
  \triangleABC
  \draw[vecteurB] (4,0)--(1,3) node[nom,right]{$\vect{BC}$};
  \draw[vecteurB] (0,0)--(-3,3) node[nom,below left]{$\vect{BC}$};
  \pt{-3,3}{above}{M}
\end{tikzpicture}}
\cas{2)}
\cotefig{%
\explique{$\vect{CB}-\vect{BA}=\vect{CB}+\vect{AB}$.}{l'opposé de $\vect{BA}$ est
$\vect{AB}$}
\explique{On trace $\vect{CB}$, puis un représentant de $\vect{AB}$ à la suite, depuis $B$
: $4$ à droite. Il arrive en $M$.}{l'un à la suite de l'autre}}
{\begin{tikzpicture}[scale=0.4]
  \quadrillage{-0.4}{-0.4}{8.4}{3.4}
  \triangleABC
  \draw[vecteurC] (1,3)--(4,0) node[nom,left]{$\vect{CB}$};
  \draw[vecteur] (4,0)--(8,0) node[nom,below]{$\vect{AB}$};
  \draw[cherche] (1,3)--(8,0) node[nom,above right]{$\vect{CM}$};
  \pt{8,0}{below}{M}
\end{tikzpicture}}
\cas{3)}
\cotefig{%
\explique{$-\vect{CB}=\vect{BC}$ ; on range les deux premiers vecteurs :
$\vect{AC}+\vect{BA}=\vect{BA}+\vect{AC}$.}{l'ordre ne compte pas}
\explique{$\vect{BA}+\vect{AC}$ va de $B$ à $C$ : c'est $\vect{BC}$.}{bout à bout}
\explique{Donc $\vect{BM}=\vect{BC}+\vect{BC}$ : depuis $B$, on reporte $\vect{BC}$ deux
fois de suite, soit $6$ à gauche et $6$ vers le haut.}{l'un à la suite de l'autre}}
{\begin{tikzpicture}[scale=0.4]
  \quadrillage{-2.4}{-0.4}{4.4}{6.4}
  \triangleABC
  \draw[cherche] (4,0)--(-2,6) node[nom,above right,pos=0.7]{$\vect{BM}$};
  \draw[vecteurB] (4,0)--(1,3) node[nom,below left]{$\vect{BC}$};
  \draw[vecteurB] (1,3)--(-2,6) node[nom,below left]{$\vect{BC}$};
  \pt{-2,6}{left}{M}
\end{tikzpicture}}
\rem{En 3), $\vect{BM}$ est caché sous les deux représentants de $\vect{BC}$ : il va de $B$
à $M$. Erreur fréquente : oublier de changer $-\vect{CB}$ en $\vect{BC}$ avant de
construire.}
\end{exo}

\begin{exo}[Symétrie centrale]
\cotefig[0.32]{%
\explique{$E'$ est le symétrique de $E$ par rapport à $G$ : $G$ est le milieu de $[EE']$.
De même, $G$ est le milieu de $[FF']$.}{définition de la symétrie centrale}}
{\begin{tikzpicture}[scale=0.45]
  \draw[fig] (2,4)--(0,1)--(4,0)--(6,3)--cycle;
  \draw[aide] (0,1)--(6,3) (2,4)--(4,0);
  \draw[vecteur] (2,4)--(0,1) node[nom,left]{$\vect{FE}$};
  \draw[vecteurB] (6,3)--(4,0) node[nom,right]{$\vect{E'F'}$};
  \pt{0,1}{left}{E} \pt{6,3}{right}{E'} \pt{2,4}{above}{F} \pt{4,0}{below}{F'}
  \pt{3,2}{above}{G}
\end{tikzpicture}}
\cas{1)}
\explique{Dans le quadrilatère $FEF'E'$, les diagonales sont $[FF']$ et $[EE']$.}{on fait
le tour $F$, $E$, $F'$, $E'$}
\explique{Elles ont le même milieu $G$.}{par définition de $E'$ et de $F'$}
\explique{Donc \res{$FEF'E'$ est un parallélogramme}.}{un quadrilatère dont les diagonales
ont le même milieu est un parallélogramme}
\cas{2) a)}
\explique{$FEF'E'$ est un parallélogramme, donc $\vect{FE}=\vect{E'F'}$.}{théorème 1 :
$ABCD$ donne $\vect{AB}=\vect{DC}$}
\cas{2) b)}
\explique{Deux vecteurs égaux ont la même norme : \res{$E'F'=EF$}.}{même longueur}
\cas{2) c)}
\explique{$[E'F']$ est l'image de $[EF]$ par la symétrie de centre $G$, et elle a la même
longueur : \res{la symétrie centrale conserve les longueurs}.}{$E$ et $F$ sont deux points
quelconques}
\rem{En 2) a), on lit bien $\vect{FE}=\vect{E'F'}$, et non $\vect{EF}=\vect{E'F'}$ : sur la
figure, $\vect{E'F'}$ est l'opposé de $\vect{EF}$. Une symétrie centrale change le sens
des vecteurs.}
\end{exo}

\partie{Multiplier un vecteur par un réel}

\begin{exo}[Multiples d'un vecteur]
\explique{$\vec{a}$ : $1$ carreau à droite et $1$ vers le bas. On multiplie les deux nombres
de carreaux par $k$ ; si $k<0$, on change le sens.}{méthode 6}
\cotefig[0.4]{%
\explique{$2\vec{a}$ : $2$ à droite et $2$ vers le bas.}{même sens, deux fois plus long}
\explique{$-\vec{a}$ : $1$ à gauche et $1$ vers le haut.}{sens contraire, même longueur}
\explique{$3\vec{a}$ : $3$ à droite et $3$ vers le bas.}{même sens, trois fois plus long}
\explique{$-2\vec{a}$ : $2$ à gauche et $2$ vers le haut.}{sens contraire, deux fois plus
long}}
{\begin{tikzpicture}[scale=0.4]
  \quadrillage{-0.4}{-0.4}{7.4}{5.4}
  \draw[cherche] (3,3)--(5,1) node[nom,below left]{$2\vec{a}$};
  \draw[vecteur] (3,3)--(4,2) node[nom,above right]{$\vec{a}$};
  \draw[cherche] (3,3)--(2,4) node[nom,below left]{$-\vec{a}$};
  \pt{3,3}{above right}{O}
\end{tikzpicture}\hspace{4pt}%
\begin{tikzpicture}[scale=0.4]
  \quadrillage{-0.4}{-0.4}{7.4}{5.4}
  \draw[cherche] (3,3)--(6,0) node[nom,below left]{$3\vec{a}$};
  \draw[vecteur] (3,3)--(4,2) node[nom,above right]{$\vec{a}$};
  \draw[cherche] (3,3)--(1,5) node[nom,below left]{$-2\vec{a}$};
  \pt{3,3}{above right}{O}
\end{tikzpicture}}
\rem{Tous ces vecteurs sont sur une même droite, celle de $\vec{a}$ : ils ont la même
direction. Deux figures pour que les flèches ne se cachent pas.}
\end{exo}

\begin{exo}[Avec des coefficients]
\explique{On trace un triangle $ABC$ sur le quadrillage : $\vect{AC}$ fait $2$ à droite et
$2$ vers le haut ; $\vect{CB}$, $2$ à droite et $2$ vers le bas ; $\vect{BA}$, $4$ à
gauche.}{des nombres de carreaux pairs, pour les moitiés}
\cas{Le point $M$}
\cotefig[0.3]{%
\explique{$3\vect{AC}$ : $6$ à droite et $6$ vers le haut, depuis $A$.}{même sens, trois
fois plus long}
\explique{$M$ est sur la demi-droite $[AC)$, avec $AM=3AC$.}{}}
{\begin{tikzpicture}[scale=0.32]
  \quadrillage{-0.4}{-0.4}{6.4}{6.4}
  \draw[fig] (0,0)--(4,0)--(2,2)--cycle;
  \draw[cherche] (0,0)--(6,6) node[nom,above left]{$3\vect{AC}$};
  \draw[vecteur] (0,0)--(2,2) node[nom,below right]{$\vect{AC}$};
  \pt{0,0}{below left}{A} \pt{4,0}{below}{B} \pt{2,2}{right}{C} \pt{6,6}{above}{M}
\end{tikzpicture}}
\cas{Le point $N$}
\cotefig[0.3]{%
\explique{$\tfrac12\vect{CB}$ : $1$ à droite et $1$ vers le bas, depuis $C$.}{même sens,
deux fois plus court}
\explique{$N$ est le milieu de $[CB]$.}{$\vect{CN}=\tfrac12\vect{CB}$}}
{\begin{tikzpicture}[scale=0.55]
  \quadrillage{-0.4}{-0.4}{4.4}{2.4}
  \draw[fig] (0,0)--(4,0)--(2,2)--cycle;
  \draw[vecteurC] (2,2)--(4,0) node[nom,above right]{$\vect{CB}$};
  \draw[cherche] (2,2)--(3,1) node[nom,below left]{$\tfrac12\vect{CB}$};
  \pt{0,0}{below left}{A} \pt{4,0}{below}{B} \pt{2,2}{above}{C} \pt{3,1}{right}{N}
\end{tikzpicture}}
\cas{Le point $P$}
\cotefig[0.3]{%
\explique{$\vect{BA}$ fait $4$ carreaux vers la gauche ; $-\tfrac12\vect{BA}$ fait donc $2$
carreaux vers la droite, depuis $B$.}{sens contraire, deux fois plus court}
\explique{$P$ est sur la droite $(AB)$, au-delà de $B$, avec $BP=\tfrac12AB$.}{}}
{\begin{tikzpicture}[scale=0.45]
  \quadrillage{-0.4}{-0.4}{6.4}{2.4}
  \draw[fig] (0,0)--(4,0)--(2,2)--cycle;
  \draw[vecteur] (4,0)--(0,0) node[nom,above]{$\vect{BA}$};
  \draw[cherche] (4,0)--(6,0) node[nom,above]{$-\tfrac12\vect{BA}$};
  \pt{0,0}{below left}{A} \pt{4,0}{below}{B} \pt{2,2}{above}{C} \pt{6,0}{below}{P}
\end{tikzpicture}}
\rem{Sans quadrillage, on mesure : pour $N$, on prend le milieu de $[CB]$ ; pour $P$, on
reporte au compas la moitié de $AB$ au-delà de $B$. Erreur fréquente pour $P$ : aller vers
$A$, en oubliant que le signe moins change le sens de $\vect{BA}$.}
\end{exo}

\begin{exo}[Sur un quadrillage]
\explique{On lit les déplacements : $\vect{BC}$, $2$ à droite et $2$ vers le haut ;
$\vect{AC}$, $6$ à droite et $2$ vers le haut ; $\vect{BD}$, $3$ à gauche et $3$ vers le
haut ; $\vect{AB}$, $4$ à droite.}{méthode 6 : on multiplie les nombres de carreaux}
\cas{1)}
\cotefig[0.36]{%
\explique{$2\vect{BC}$ : $4$ à droite et $4$ vers le haut, depuis $A$.}{même sens, deux
fois plus long}}
{\begin{tikzpicture}[scale=0.36]\figABCD
  \draw[vecteurB] (4,1)--(6,3) node[nom,right]{$\vect{BC}$};
  \draw[cherche] (0,1)--(4,5) node[nom,above left]{$2\vect{BC}$};
\end{tikzpicture}}
\cas{2)}
\cotefig[0.36]{%
\explique{$0{,}5\,\vect{AC}$ : $3$ à droite et $1$ vers le haut, depuis $B$.}{même sens,
deux fois plus court}}
{\begin{tikzpicture}[scale=0.36]\figABCD
  \draw[vecteur] (0,1)--(6,3) node[nom,above left]{$\vect{AC}$};
  \draw[cherche] (4,1)--(7,2) node[nom,below right]{$0{,}5\,\vect{AC}$};
\end{tikzpicture}}
\cas{3)}
\cotefig[0.36]{%
\explique{$-\tfrac13\vect{BD}$ : $1$ à droite et $1$ vers le bas, depuis $D$.}{sens
contraire, trois fois plus court}}
{\begin{tikzpicture}[scale=0.36]\figABCD
  \draw[vecteurC] (4,1)--(1,4) node[nom,below left]{$\vect{BD}$};
  \draw[cherche] (1,4)--(2,3) node[nom,above right]{$-\tfrac13\vect{BD}$};
\end{tikzpicture}}
\cas{4)}
\cotefig[0.36]{%
\explique{$\tfrac13\vect{BD}$ : $1$ à gauche et $1$ vers le haut, depuis $C$.}{même sens,
trois fois plus court}}
{\begin{tikzpicture}[scale=0.36]\figABCD
  \draw[vecteurC] (4,1)--(1,4) node[nom,below left]{$\vect{BD}$};
  \draw[cherche] (6,3)--(5,4) node[nom,above right]{$\tfrac13\vect{BD}$};
\end{tikzpicture}}
\cas{5)}
\cotefig[0.36]{%
\explique{$\vect{AB}-\vect{BC}=\vect{AB}+(-\vect{BC})$ : depuis $D$, on trace $\vect{AB}$
($4$ à droite), puis $-\vect{BC}$ à la suite ($2$ à gauche et $2$ vers le bas).}{méthode 4
: l'opposé à la suite}
\explique{Le vecteur cherché va de $D$ à l'arrivée : $2$ à droite et $2$ vers le
bas.}{départ, arrivée}}
{\begin{tikzpicture}[scale=0.36]\figABCD
  \draw[vecteur] (1,4)--(5,4) node[nom,above]{$\vect{AB}$};
  \draw[vecteurB] (5,4)--(3,2) node[nom,right]{$-\vect{BC}$};
  \draw[cherche] (1,4)--(3,2) node[nom,below left]{$\vect{AB}-\vect{BC}$};
\end{tikzpicture}}
\rem{Le signe moins ne change que le sens : il ne raccourcit pas le vecteur. On vérifie que
chaque extrémité tombe sur un nœud du quadrillage, comme l'annonce l'énoncé.}
\end{exo}

\partie{Relation de Chasles}

\begin{exo}[Un seul vecteur, si c'est possible]
\explique{La relation de Chasles s'applique quand la lettre d'arrivée du premier vecteur
est la lettre de départ du second ; on peut d'abord échanger l'ordre.}{théorème 3 :
$\vect{AB}+\vect{BC}=\vect{AC}$}
\cas{a)}
\begin{calculs}
\vect{BE}+\vect{EG} &= \res{\vect{BG}} &&\justif{déjà rangés : $B\to E\to G$}
\end{calculs}
\cas{b)}
\explique{$\vect{DA}+\vect{DC}$ : \res{impossible}.}{les deux vecteurs partent de $D$ :
aucune lettre ne relie l'arrivée de l'un au départ de l'autre}
\cas{c)}
\begin{calculs}
\vect{TU}+\vect{KT} &= \vect{KT}+\vect{TU} &&\justif{on range : $K\to T\to U$}\\
\vect{TU}+\vect{KT} &= \res{\vect{KU}} &&\justif{Chasles}
\end{calculs}
\cas{d)}
\begin{calculs}
\vect{FH}+\vect{HS} &= \res{\vect{FS}} &&\justif{déjà rangés : $F\to H\to S$}
\end{calculs}
\cas{e)}
\begin{calculs}
\vect{PQ}+\vect{QP} &= \vect{PP} &&\justif{aller de $P$ à $Q$, puis revenir en $P$}\\
\vect{PQ}+\vect{QP} &= \res{\vec{0}} &&\justif{on est revenu au point de départ}
\end{calculs}
\cas{f)}
\explique{$\vect{AB}+\vect{CE}$ : \res{impossible}.}{l'arrivée $B$ n'est pas le départ $C$,
et $E$ n'est pas $A$}
\rem{En b), la somme existe bien (méthode 4), mais elle ne s'écrit pas avec les lettres de
l'énoncé : il faudrait construire un nouveau point.}
\end{exo}

\begin{exo}[Simplifier]
\explique{On remplace chaque $-\vect{XY}$ par $\vect{YX}$, puis on range les vecteurs pour
que l'extrémité de chacun soit l'origine du suivant.}{méthode 7}
\cas{$\vec{u}$}
\begin{calculs}
\vec{u} &= \vect{BC}+\vect{CA}+\vect{AB} &&\justif{déjà rangés : $B\to C\to A\to B$}\\
\vec{u} &= \vect{BB} &&\justif{on revient au point de départ}\\
\vec{u} &= \res{\vec{0}} &&
\end{calculs}
\cas{$\vec{v}$}
\begin{calculs}
\vec{v} &= \vect{CA}-\vect{CB}+\vect{AC} &&\justif{on recopie}\\
\vec{v} &= \vect{CA}+\vect{BC}+\vect{AC} &&\justif{$-\vect{CB}=\vect{BC}$}\\
\vec{v} &= \vect{BC}+\vect{CA}+\vect{AC} &&\justif{on range : $B\to C\to A\to C$}\\
\vec{v} &= \vect{BA}+\vect{AC} &&\justif{Chasles : $\vect{BC}+\vect{CA}=\vect{BA}$}\\
\vec{v} &= \res{\vect{BC}} &&\justif{Chasles}
\end{calculs}
\cas{$\vec{w}$}
\begin{calculs}
\vec{w} &= \vect{MC}-\vect{MA}+\vect{CB} &&\justif{on recopie}\\
\vec{w} &= \vect{MC}+\vect{AM}+\vect{CB} &&\justif{$-\vect{MA}=\vect{AM}$}\\
\vec{w} &= \vect{AM}+\vect{MC}+\vect{CB} &&\justif{on range : $A\to M\to C\to B$}\\
\vec{w} &= \res{\vect{AB}} &&\justif{Chasles, deux fois}
\end{calculs}
\rem{Le point $M$ disparaît : on passe par $M$ pour aller de $A$ à $C$. Erreur fréquente :
garder $-\vect{MA}$ et tenter Chasles avec un signe moins ; on change d'abord le signe en
inversant les lettres.}
\end{exo}

\begin{exo}[Compléter]
\explique{On cherche la lettre qui fait un chemin : l'arrivée d'un vecteur est le départ du
suivant.}{Chasles : $\vect{AB}+\vect{BC}=\vect{AC}$}
\cas{1)}
\explique{$\vect{BC}=\vect{BA}+\vect{A\res{C}}$.}{de $B$ à $A$, puis de $A$ à $C$ : on va de
$B$ à $C$}
\cas{2)}
\explique{$\vect{CB}+\vect{\res{B}A}=\vect{CA}$.}{le second vecteur part de $B$, où
arrive le premier}
\cas{3)}
\explique{$\vect{A\res{B}}+\vect{CA}=\vect{CB}$.}{rangé : $\vect{CA}+\vect{AB}=\vect{CB}$}
\cas{4)}
\begin{calculs}
\vect{BC}+\vect{DA}+\vect{CD} &= \vect{BC}+\vect{CD}+\vect{DA} &&\justif{on range : $B\to C\to D\to A$}\\
\vect{BC}+\vect{DA}+\vect{CD} &= \res{\vect{BA}} &&\justif{Chasles}
\end{calculs}
\rem{En 3), on commence par ranger : le vecteur $\vect{CA}$ part de $C$, comme le résultat
$\vect{CB}$ ; il vient donc en premier.}
\end{exo}

\begin{exo}[Peut-on réduire ?]
\cas{a)}
\begin{calculs}
\vect{MN}+\vect{NP} &= \res{\vect{MP}} &&\justif{oui : $M\to N\to P$}
\end{calculs}
\cas{b)}
\explique{$\vect{MN}+\vect{MP}$ : \res{non}.}{même origine $M$ : la somme suit la diagonale
d'un parallélogramme, mais ne s'écrit pas avec ces lettres}
\cas{c)}
\begin{calculs}
\vect{RO}+\vect{OS} &= \res{\vect{RS}} &&\justif{oui : $R\to O\to S$}
\end{calculs}
\cas{d)}
\begin{calculs}
\vect{RO}+\vect{SR} &= \vect{SR}+\vect{RO} &&\justif{oui, après avoir rangé : $S\to R\to O$}\\
\vect{RO}+\vect{SR} &= \res{\vect{SO}} &&\justif{Chasles}
\end{calculs}
\cas{e)}
\begin{calculs}
\vect{GH}+\vect{HG} &= \vect{GG} &&\justif{oui : aller-retour}\\
\vect{GH}+\vect{HG} &= \res{\vec{0}} &&
\end{calculs}
\cas{f)}
\begin{calculs}
\vect{KB}+\vect{AK} &= \vect{AK}+\vect{KB} &&\justif{oui, après avoir rangé : $A\to K\to B$}\\
\vect{KB}+\vect{AK} &= \res{\vect{AB}} &&\justif{Chasles}
\end{calculs}
\rem{Avant de répondre « non », on essaie toujours d'échanger les deux vecteurs (d et f).}
\end{exo}

\begin{exo}[Une égalité à démontrer]
\explique{On part du membre de gauche et on le transforme jusqu'à obtenir celui de
droite.}{on ne part pas du résultat}
\begin{calculs}
\vect{OB}-\vect{OA}+\vect{BC} &= \vect{OB}+\vect{AO}+\vect{BC} &&\justif{$-\vect{OA}=\vect{AO}$}\\
\vect{OB}-\vect{OA}+\vect{BC} &= \vect{AO}+\vect{OB}+\vect{BC} &&\justif{on range : $A\to O\to B\to C$}\\
\vect{OB}-\vect{OA}+\vect{BC} &= \vect{AC} &&\justif{Chasles, deux fois}
\end{calculs}
\explique{L'égalité \res{$\vect{OB}-\vect{OA}+\vect{BC}=\vect{AC}$} est démontrée, pour tous
points $O$, $A$, $B$, $C$.}{aucune figure particulière n'a servi}
\end{exo}

\begin{exo}[Exprimer avec deux vecteurs]
\explique{On fait apparaître les vecteurs demandés avec Chasles (on passe par le point
utile), puis on calcule comme en calcul littéral.}{$k(\vec{u}+\vec{v})=k\vec{u}+k\vec{v}$}
\cas{1) a)}
\begin{calculs}
\vec{u} &= \vect{CB} &&\justif{on recopie}\\
\vec{u} &= \vect{CA}+\vect{AB} &&\justif{Chasles : on passe par $A$}\\
\vec{u} &= -\vect{AC}+\vect{AB} &&\justif{$\vect{CA}=-\vect{AC}$}\\
\vec{u} &= \res{\vect{AB}-\vect{AC}} &&
\end{calculs}
\cas{1) b)}
\begin{calculs}
\vec{u} &= 3\vect{CB}+\vect{BA} &&\justif{on recopie}\\
\vec{u} &= 3\bigl(\vect{AB}-\vect{AC}\bigr)-\vect{AB} &&\justif{d'après a) ; $\vect{BA}=-\vect{AB}$}\\
\vec{u} &= 3\vect{AB}-3\vect{AC}-\vect{AB} &&\justif{on développe}\\
\vec{u} &= \res{2\vect{AB}-3\vect{AC}} &&\justif{on réduit : $3-1=2$}
\end{calculs}
\cas{2) a)}
\begin{calculs}
\vec{v} &= \vect{BA}+\vect{CA} &&\justif{on recopie}\\
\vec{v} &= \vect{BC}+\vect{CA}+\vect{CA} &&\justif{Chasles : $\vect{BA}=\vect{BC}+\vect{CA}$}\\
\vec{v} &= \res{\vect{BC}+2\vect{CA}} &&\justif{on réduit}
\end{calculs}
\cas{2) b)}
\begin{calculs}
\vec{v} &= \vect{AB}+2\vect{BC}+\vect{AC} &&\justif{on recopie}\\
\vec{v} &= \vect{AB}+\vect{BC}+\vect{BC}+\vect{AC} &&\justif{$2\vect{BC}=\vect{BC}+\vect{BC}$}\\
\vec{v} &= \vect{AC}+\vect{BC}+\vect{AC} &&\justif{Chasles : $\vect{AB}+\vect{BC}=\vect{AC}$}\\
\vec{v} &= \vect{BC}+2\vect{AC} &&\justif{on réduit}\\
\vec{v} &= \res{\vect{BC}-2\vect{CA}} &&\justif{$\vect{AC}=-\vect{CA}$}
\end{calculs}
\rem{On repère d'abord les deux vecteurs autorisés, puis on passe par le point commun : en
1), tout part de $A$ ; en 2), on passe par $C$. On peut vérifier 1) b) sur une figure en
construisant les deux membres depuis un même point.}
\end{exo}

\begin{exo}[Différences sur un quadrillage]
\explique{Pour chaque différence, on change le signe en inversant les lettres, puis on
range avec Chasles quand c'est possible : le vecteur cherché devient un vecteur de la
figure, qu'on reporte depuis $A$.}{méthode 7, puis méthode 2}
\cas{a)}
\cotefig[0.36]{%
\begin{calculs}
\vect{AC}-\vect{AB} &= \vect{AC}+\vect{BA} &&\justif{$-\vect{AB}=\vect{BA}$}\\
\vect{AC}-\vect{AB} &= \vect{BA}+\vect{AC} &&\justif{on range}\\
\vect{AC}-\vect{AB} &= \vect{BC} &&\justif{Chasles}
\end{calculs}
\explique{$\vect{BC}$ : $1$ à droite et $3$ vers le haut ; on le reporte depuis
$A$.}{}}
{\begin{tikzpicture}[scale=0.34]\figABCDbis
  \draw[vecteurB] (4,0)--(5,3) node[nom,right]{$\vect{BC}$};
  \draw[cherche] (0,1)--(1,4) node[nom,left]{$\vect{BC}$};
\end{tikzpicture}}
\cas{b)}
\cotefig[0.36]{%
\begin{calculs}
\vect{AD}-\vect{BD} &= \vect{AD}+\vect{DB} &&\justif{$-\vect{BD}=\vect{DB}$}\\
\vect{AD}-\vect{BD} &= \vect{AB} &&\justif{Chasles : $A\to D\to B$}
\end{calculs}
\explique{Le représentant d'origine $A$ est $\vect{AB}$ lui-même : la flèche arrive en
$B$.}{rien à reporter}}
{\begin{tikzpicture}[scale=0.34]\figABCDbis
  \draw[cherche] (0,1)--(4,0) node[nom,below left]{$\vect{AB}$};
\end{tikzpicture}}
\cas{c)}
\cotefig[0.36]{%
\begin{calculs}
\vect{AB}-\vect{AD} &= \vect{AB}+\vect{DA} &&\justif{$-\vect{AD}=\vect{DA}$}\\
\vect{AB}-\vect{AD} &= \vect{DA}+\vect{AB} &&\justif{on range}\\
\vect{AB}-\vect{AD} &= \vect{DB} &&\justif{Chasles}
\end{calculs}
\explique{$\vect{DB}$ : $3$ à droite et $3$ vers le bas ; on le reporte depuis
$A$.}{on sort du quadrillage : on le prolonge}}
{\begin{tikzpicture}[scale=0.34]\figABCDbis
  \draw[vecteurC] (1,3)--(4,0) node[nom,above right]{$\vect{DB}$};
  \draw[cherche] (0,1)--(3,-2) node[nom,below left]{$\vect{DB}$};
\end{tikzpicture}}
\cas{d)}
\cotefig[0.36]{%
\begin{calculs}
\vect{AC}-\vect{BD} &= \vect{AC}+\vect{DB} &&\justif{$-\vect{BD}=\vect{DB}$}
\end{calculs}
\explique{Chasles ne s'applique pas : $C$ n'est pas $D$. On construit : depuis $C$, on trace
$\vect{DB}$ à la suite de $\vect{AC}$.}{méthode 4 : l'un à la suite de l'autre}
\explique{Le vecteur cherché va de $A$ à l'arrivée : $8$ à droite et $1$ vers le
bas.}{départ, arrivée}}
{\begin{tikzpicture}[scale=0.34]\figABCDbis
  \draw[vecteur] (0,1)--(5,3) node[nom,above left]{$\vect{AC}$};
  \draw[vecteurC] (5,3)--(8,0) node[nom,above right]{$\vect{DB}$};
  \draw[cherche] (0,1)--(8,0) node[nom,below,pos=0.75]{$\vect{AC}-\vect{BD}$};
\end{tikzpicture}}
\rem{Erreur fréquente en b) : chercher à tracer quelque chose, alors que le calcul donne
directement un vecteur d'origine $A$. Le calcul d'abord, la construction ensuite.}
\end{exo}

\end{document}
